%%%%%%%%%%%%%%%%%%%%%%%%%%%%%%%%%%%%%%%%%%%%%%%%%%%%%%%%%%%%
% 12.1 FORMAL LOGIC
% The Composition of Advanced Mathematics
%%%%%%%%%%%%%%%%%%%%%%%%%%%%%%%%%%%%%%%%%%%%%%%%%%%%%%%%%%%%

\section{Formal Logic}
\label{sec:formal-logic}

\prerequisite{
Elementary logic, sets, functions, relations, methods of proof, abstract
algebraic notation, and familiarity with mathematical definitions and
theorems.
}

\learningobjectives{
By the end of this section, the reader should be able to:
\begin{itemize}
    \item distinguish formal logic from informal mathematical reasoning;
    \item distinguish syntax from semantics;
    \item identify formal languages and alphabets;
    \item construct well-formed formulas;
    \item distinguish object language and metalanguage;
    \item understand propositional variables and connectives;
    \item construct truth assignments;
    \item define semantic satisfaction;
    \item distinguish satisfiable, unsatisfiable, and valid formulas;
    \item understand logical consequence;
    \item recognize logical equivalence;
    \item derive standard propositional equivalences;
    \item understand normal forms;
    \item construct conjunctive and disjunctive normal forms;
    \item recognize Boolean functions;
    \item understand functionally complete connective sets;
    \item recognize the satisfiability problem;
    \item understand formal proof systems;
    \item distinguish axioms, inference rules, and derivations;
    \item understand syntactic derivability;
    \item distinguish derivability from semantic consequence;
    \item recognize soundness and completeness;
    \item understand consistency conceptually;
    \item recognize deduction theorems;
    \item understand natural deduction;
    \item recognize introduction and elimination rules;
    \item understand proof by contradiction formally;
    \item understand first-order languages;
    \item distinguish variables, constants, function symbols, and relation
          symbols;
    \item construct terms and formulas recursively;
    \item distinguish free and bound variables;
    \item understand substitution;
    \item recognize variable capture;
    \item understand first-order structures and interpretations;
    \item evaluate terms in structures;
    \item define satisfaction for first-order formulas;
    \item understand quantifier semantics;
    \item distinguish sentences from open formulas;
    \item understand models and theories;
    \item recognize valid and satisfiable first-order formulas;
    \item understand semantic consequence for theories;
    \item recognize first-order proof systems;
    \item understand first-order soundness and completeness conceptually;
    \item recognize compactness conceptually;
    \item recognize Löwenheim--Skolem phenomena conceptually;
    \item understand equality in first-order logic;
    \item translate mathematical statements into formal language;
    \item recognize expressive limitations of first-order logic;
    \item connect formal logic with set theory, model theory, proof theory,
          computability, type theory, and foundations of mathematics.
\end{itemize}
}

%%%%%%%%%%%%%%%%%%%%%%%%%%%%%%%%%%%%%%%%%%%%%%%%%%%%%%%%%%%%
\subsection{From Elementary Logic to Formal Logic}
\label{subsec:elementary-to-formal-logic}
%%%%%%%%%%%%%%%%%%%%%%%%%%%%%%%%%%%%%%%%%%%%%%%%%%%%%%%%%%%%

Elementary logic studies common reasoning patterns such as

\[
P\land Q,
\qquad
P\lor Q,
\qquad
P\Rightarrow Q,
\]

and quantified statements such as

\[
\forall x\,P(x).
\]

Formal logic goes further by treating the language of mathematics itself
as a mathematical object.

Instead of merely reasoning with formulas, we study:

\[
\boxed{
\text{how formulas are constructed},
}
\]

\[
\boxed{
\text{what formulas mean},
}
\]

and

\[
\boxed{
\text{how proofs are generated}.
}
\]

%%%%%%%%%%%%%%%%%%%%%%%%%%%%%%%%%%%%%%%%%%%%%%%%%%%%%%%%%%%%
\subsection{Syntax and Semantics}
\label{subsec:syntax-and-semantics}
%%%%%%%%%%%%%%%%%%%%%%%%%%%%%%%%%%%%%%%%%%%%%%%%%%%%%%%%%%%%

Formal logic separates two fundamental levels.

\begin{definitionbox}{Syntax}
Syntax is the formal structure of expressions and proofs.

It describes which sequences of symbols are legal formulas and which
sequences of formulas form valid derivations.
\end{definitionbox}

\begin{definitionbox}{Semantics}
Semantics assigns mathematical meaning or truth conditions to formal
expressions.
\end{definitionbox}

Thus

\[
\boxed{
\text{syntax}
=
\text{formal symbols and rules},
}
\]

while

\[
\boxed{
\text{semantics}
=
\text{interpretation and truth}.
}
\]

%%%%%%%%%%%%%%%%%%%%%%%%%%%%%%%%%%%%%%%%%%%%%%%%%%%%%%%%%%%%
\subsection{Object Language and Metalanguage}
\label{subsec:object-language-metalanguage}
%%%%%%%%%%%%%%%%%%%%%%%%%%%%%%%%%%%%%%%%%%%%%%%%%%%%%%%%%%%%

The \textbf{object language} is the formal language being studied.

The \textbf{metalanguage} is the language used to discuss the object
language.

For example,

\[
P\land Q
\]

may be a formula of the object language.

The statement

\[
\boxed{
\text{``\(P\land Q\) is satisfiable''}
}
\]

belongs to the metalanguage.

Confusing these two levels can create logical ambiguity.

%%%%%%%%%%%%%%%%%%%%%%%%%%%%%%%%%%%%%%%%%%%%%%%%%%%%%%%%%%%%
\subsection{Formal Languages}
\label{subsec:formal-languages}
%%%%%%%%%%%%%%%%%%%%%%%%%%%%%%%%%%%%%%%%%%%%%%%%%%%%%%%%%%%%

A formal language begins with an alphabet of symbols.

The alphabet may include:

\[
\boxed{
\text{variables},
}
\]

\[
\boxed{
\text{logical connectives},
}
\]

\[
\boxed{
\text{quantifiers},
}
\]

\[
\boxed{
\text{parentheses},
}
\]

and, in first-order logic,

\[
\boxed{
\text{function and relation symbols}.
}
\]

Rules then determine which finite strings are well-formed.

%%%%%%%%%%%%%%%%%%%%%%%%%%%%%%%%%%%%%%%%%%%%%%%%%%%%%%%%%%%%
\subsection{Recursive Definitions}
\label{subsec:formal-recursive-definitions}
%%%%%%%%%%%%%%%%%%%%%%%%%%%%%%%%%%%%%%%%%%%%%%%%%%%%%%%%%%%%

Logical syntax is usually defined recursively.

A recursive definition contains:

\begin{enumerate}
    \item basic objects;
    \item rules for constructing more complicated objects;
    \item a statement that nothing else is included.
\end{enumerate}

This permits structural induction and recursive algorithms on formulas.

%%%%%%%%%%%%%%%%%%%%%%%%%%%%%%%%%%%%%%%%%%%%%%%%%%%%%%%%%%%%
\subsection{Propositional Logic}
\label{subsec:formal-propositional-logic}
%%%%%%%%%%%%%%%%%%%%%%%%%%%%%%%%%%%%%%%%%%%%%%%%%%%%%%%%%%%%

Propositional logic begins with propositional variables

\[
\boxed{
P,Q,R,\ldots
}
\]

whose truth values are

\[
\boxed{
\text{True}
\quad\text{or}\quad
\text{False}.
}
\]

Complex formulas are formed using logical connectives.

%%%%%%%%%%%%%%%%%%%%%%%%%%%%%%%%%%%%%%%%%%%%%%%%%%%%%%%%%%%%
\subsection{Logical Connectives}
\label{subsec:formal-logical-connectives}
%%%%%%%%%%%%%%%%%%%%%%%%%%%%%%%%%%%%%%%%%%%%%%%%%%%%%%%%%%%%

The standard connectives are:

\[
\boxed{
\neg P
}
\]

for negation,

\[
\boxed{
P\land Q
}
\]

for conjunction,

\[
\boxed{
P\lor Q
}
\]

for disjunction,

\[
\boxed{
P\rightarrow Q
}
\]

for implication, and

\[
\boxed{
P\leftrightarrow Q
}
\]

for biconditional.

%%%%%%%%%%%%%%%%%%%%%%%%%%%%%%%%%%%%%%%%%%%%%%%%%%%%%%%%%%%%
\subsection{Well-Formed Propositional Formulas}
\label{subsec:wff-propositional}
%%%%%%%%%%%%%%%%%%%%%%%%%%%%%%%%%%%%%%%%%%%%%%%%%%%%%%%%%%%%

A propositional formula is defined recursively:

\begin{enumerate}
    \item every propositional variable is a formula;
    \item if \(\varphi\) is a formula, then
    \[
    \neg\varphi
    \]
    is a formula;
    \item if \(\varphi\) and \(\psi\) are formulas, then
    \[
    (\varphi\land\psi),
    \quad
    (\varphi\lor\psi),
    \quad
    (\varphi\rightarrow\psi),
    \quad
    (\varphi\leftrightarrow\psi)
    \]
    are formulas;
    \item no other strings are formulas.
\end{enumerate}

Such formulas are often called \textbf{well-formed formulas}, abbreviated
WFFs.

%%%%%%%%%%%%%%%%%%%%%%%%%%%%%%%%%%%%%%%%%%%%%%%%%%%%%%%%%%%%
\subsection{Parse Trees}
\label{subsec:logical-parse-trees}
%%%%%%%%%%%%%%%%%%%%%%%%%%%%%%%%%%%%%%%%%%%%%%%%%%%%%%%%%%%%

A formula has hierarchical structure.

For

\[
\neg(P\land(Q\rightarrow R)),
\]

the outermost connective is

\[
\boxed{
\neg.
}
\]

Its immediate subformula is

\[
P\land(Q\rightarrow R).
\]

That formula has two major components:

\[
P
\]

and

\[
Q\rightarrow R.
\]

A parse tree records this recursive structure.

%%%%%%%%%%%%%%%%%%%%%%%%%%%%%%%%%%%%%%%%%%%%%%%%%%%%%%%%%%%%
\subsection{Unique Readability}
\label{subsec:unique-readability}
%%%%%%%%%%%%%%%%%%%%%%%%%%%%%%%%%%%%%%%%%%%%%%%%%%%%%%%%%%%%

A carefully defined formal syntax satisfies a unique-readability property:

a well-formed formula can be parsed in exactly one permitted way.

For example,

\[
(P\land Q)\lor R
\]

and

\[
P\land(Q\lor R)
\]

are different formulas.

Parentheses or precedence conventions prevent ambiguity.

%%%%%%%%%%%%%%%%%%%%%%%%%%%%%%%%%%%%%%%%%%%%%%%%%%%%%%%%%%%%
\subsection{Subformulas}
\label{subsec:subformulas}
%%%%%%%%%%%%%%%%%%%%%%%%%%%%%%%%%%%%%%%%%%%%%%%%%%%%%%%%%%%%

A subformula is a formula occurring as a constituent of another formula.

For

\[
\varphi
=
(P\rightarrow Q)\land\neg R,
\]

subformulas include

\[
P,
\qquad
Q,
\qquad
P\rightarrow Q,
\qquad
R,
\qquad
\neg R,
\]

and the entire formula \(\varphi\).

Subformula structure is important in proof theory and satisfiability
algorithms.

%%%%%%%%%%%%%%%%%%%%%%%%%%%%%%%%%%%%%%%%%%%%%%%%%%%%%%%%%%%%
\subsection{Truth Assignments}
\label{subsec:truth-assignments}
%%%%%%%%%%%%%%%%%%%%%%%%%%%%%%%%%%%%%%%%%%%%%%%%%%%%%%%%%%%%

A truth assignment is a function

\[
\boxed{
v:
\mathcal{P}
\rightarrow
\{0,1\},
}
\]

where \(\mathcal{P}\) is the set of propositional variables.

Interpret

\[
1
=
\text{True},
\]

and

\[
0
=
\text{False}.
\]

The assignment is extended recursively to all formulas.

%%%%%%%%%%%%%%%%%%%%%%%%%%%%%%%%%%%%%%%%%%%%%%%%%%%%%%%%%%%%
\subsection{Semantic Evaluation}
\label{subsec:propositional-semantic-evaluation}
%%%%%%%%%%%%%%%%%%%%%%%%%%%%%%%%%%%%%%%%%%%%%%%%%%%%%%%%%%%%

For negation,

\[
\boxed{
v(\neg\varphi)
=
1-v(\varphi).
}
\]

For conjunction,

\[
\boxed{
v(\varphi\land\psi)=1
}
\]

exactly when

\[
v(\varphi)=v(\psi)=1.
\]

For disjunction,

\[
\boxed{
v(\varphi\lor\psi)=1
}
\]

when at least one of the two formulas is true.

%%%%%%%%%%%%%%%%%%%%%%%%%%%%%%%%%%%%%%%%%%%%%%%%%%%%%%%%%%%%
\subsection{Implication Semantics}
\label{subsec:formal-implication-semantics}
%%%%%%%%%%%%%%%%%%%%%%%%%%%%%%%%%%%%%%%%%%%%%%%%%%%%%%%%%%%%

The implication

\[
\varphi\rightarrow\psi
\]

is false only when

\[
\boxed{
v(\varphi)=1,
\qquad
v(\psi)=0.
}
\]

Equivalently,

\[
\boxed{
\varphi\rightarrow\psi
\equiv
\neg\varphi\lor\psi.
}
\]

This equivalence is central to propositional simplification.

%%%%%%%%%%%%%%%%%%%%%%%%%%%%%%%%%%%%%%%%%%%%%%%%%%%%%%%%%%%%
\subsection{Satisfaction}
\label{subsec:propositional-satisfaction}
%%%%%%%%%%%%%%%%%%%%%%%%%%%%%%%%%%%%%%%%%%%%%%%%%%%%%%%%%%%%

If truth assignment \(v\) makes formula \(\varphi\) true, write

\[
\boxed{
v\models\varphi.
}
\]

This is read:

\[
\boxed{
\text{``\(v\) satisfies \(\varphi\).''}
}
\]

The symbol

\[
\models
\]

is semantic notation.

%%%%%%%%%%%%%%%%%%%%%%%%%%%%%%%%%%%%%%%%%%%%%%%%%%%%%%%%%%%%
\subsection{Satisfiable Formulas}
\label{subsec:satisfiable-formulas}
%%%%%%%%%%%%%%%%%%%%%%%%%%%%%%%%%%%%%%%%%%%%%%%%%%%%%%%%%%%%

\begin{definitionbox}{Satisfiable Formula}
A formula \(\varphi\) is satisfiable if there exists at least one truth
assignment \(v\) such that

\[
v\models\varphi.
\]
\end{definitionbox}

For example,

\[
P\land Q
\]

is satisfiable because both \(P\) and \(Q\) can be assigned True.

%%%%%%%%%%%%%%%%%%%%%%%%%%%%%%%%%%%%%%%%%%%%%%%%%%%%%%%%%%%%
\subsection{Unsatisfiable Formulas}
\label{subsec:unsatisfiable-formulas}
%%%%%%%%%%%%%%%%%%%%%%%%%%%%%%%%%%%%%%%%%%%%%%%%%%%%%%%%%%%%

\begin{definitionbox}{Unsatisfiable Formula}
A formula \(\varphi\) is unsatisfiable if no truth assignment satisfies it.
\end{definitionbox}

For example,

\[
\boxed{
P\land\neg P
}
\]

is unsatisfiable.

Such a formula is also called a contradiction.

%%%%%%%%%%%%%%%%%%%%%%%%%%%%%%%%%%%%%%%%%%%%%%%%%%%%%%%%%%%%
\subsection{Valid Formulas}
\label{subsec:valid-formulas}
%%%%%%%%%%%%%%%%%%%%%%%%%%%%%%%%%%%%%%%%%%%%%%%%%%%%%%%%%%%%

\begin{definitionbox}{Validity}
A formula \(\varphi\) is valid if

\[
v\models\varphi
\]

for every truth assignment \(v\).
\end{definitionbox}

Write

\[
\boxed{
\models\varphi.
}
\]

A valid propositional formula is often called a tautology.

%%%%%%%%%%%%%%%%%%%%%%%%%%%%%%%%%%%%%%%%%%%%%%%%%%%%%%%%%%%%
\subsection{Worked Example: Validity}
\label{subsec:worked-validity}
%%%%%%%%%%%%%%%%%%%%%%%%%%%%%%%%%%%%%%%%%%%%%%%%%%%%%%%%%%%%

\begin{workedexamplebox}{Law of Excluded Middle}

Consider

\[
P\lor\neg P.
\]

If \(P\) is true, the first disjunct is true.

If \(P\) is false, then \(\neg P\) is true.

Thus every truth assignment satisfies the formula.

\begin{answerbox}
\[
\boxed{
\models
P\lor\neg P.
}
\]
\end{answerbox}

Therefore the formula is valid in classical propositional logic.

\end{workedexamplebox}

%%%%%%%%%%%%%%%%%%%%%%%%%%%%%%%%%%%%%%%%%%%%%%%%%%%%%%%%%%%%
\subsection{Duality Between Validity and Unsatisfiability}
\label{subsec:validity-unsatisfiability-duality}
%%%%%%%%%%%%%%%%%%%%%%%%%%%%%%%%%%%%%%%%%%%%%%%%%%%%%%%%%%%%

A formula \(\varphi\) is valid exactly when its negation is unsatisfiable:

\[
\boxed{
\models\varphi
\iff
\neg\varphi
\text{ is unsatisfiable}.
}
\]

Likewise,

\[
\boxed{
\varphi
\text{ is unsatisfiable}
\iff
\models\neg\varphi.
}
\]

This connects validity checking with satisfiability testing.

%%%%%%%%%%%%%%%%%%%%%%%%%%%%%%%%%%%%%%%%%%%%%%%%%%%%%%%%%%%%
\subsection{Logical Consequence}
\label{subsec:logical-consequence}
%%%%%%%%%%%%%%%%%%%%%%%%%%%%%%%%%%%%%%%%%%%%%%%%%%%%%%%%%%%%

Let

\[
\Gamma
\]

be a collection of formulas.

Write

\[
\boxed{
\Gamma\models\varphi
}
\]

if every truth assignment satisfying all formulas in \(\Gamma\) also
satisfies \(\varphi\).

This means \(\varphi\) is a semantic consequence of \(\Gamma\).

%%%%%%%%%%%%%%%%%%%%%%%%%%%%%%%%%%%%%%%%%%%%%%%%%%%%%%%%%%%%
\subsection{Worked Example: Semantic Consequence}
\label{subsec:worked-semantic-consequence}
%%%%%%%%%%%%%%%%%%%%%%%%%%%%%%%%%%%%%%%%%%%%%%%%%%%%%%%%%%%%

\begin{workedexamplebox}{Modus Ponens Semantically}

Let

\[
\Gamma
=
\{
P,
P\rightarrow Q
\}.
\]

Any assignment satisfying \(P\) makes \(P\) true.

For the implication to also be true, \(Q\) must be true.

Therefore,

\begin{answerbox}
\[
\boxed{
\{
P,
P\rightarrow Q
\}
\models
Q.
}
\]
\end{answerbox}

\end{workedexamplebox}

%%%%%%%%%%%%%%%%%%%%%%%%%%%%%%%%%%%%%%%%%%%%%%%%%%%%%%%%%%%%
\subsection{Logical Equivalence}
\label{subsec:formal-logical-equivalence}
%%%%%%%%%%%%%%%%%%%%%%%%%%%%%%%%%%%%%%%%%%%%%%%%%%%%%%%%%%%%

Two formulas \(\varphi\) and \(\psi\) are logically equivalent if they
have the same truth value under every assignment.

Write

\[
\boxed{
\varphi
\equiv
\psi.
}
\]

Equivalently,

\[
\boxed{
\models
\varphi\leftrightarrow\psi.
}
\]

%%%%%%%%%%%%%%%%%%%%%%%%%%%%%%%%%%%%%%%%%%%%%%%%%%%%%%%%%%%%
\subsection{Important Propositional Equivalences}
\label{subsec:propositional-equivalences}
%%%%%%%%%%%%%%%%%%%%%%%%%%%%%%%%%%%%%%%%%%%%%%%%%%%%%%%%%%%%

Double negation:

\[
\boxed{
\neg\neg P
\equiv
P.
}
\]

De Morgan's laws:

\[
\boxed{
\neg(P\land Q)
\equiv
\neg P\lor\neg Q,
}
\]

\[
\boxed{
\neg(P\lor Q)
\equiv
\neg P\land\neg Q.
}
\]

Implication:

\[
\boxed{
P\rightarrow Q
\equiv
\neg P\lor Q.
}
\]

Contraposition:

\[
\boxed{
P\rightarrow Q
\equiv
\neg Q\rightarrow\neg P.
}
\]

%%%%%%%%%%%%%%%%%%%%%%%%%%%%%%%%%%%%%%%%%%%%%%%%%%%%%%%%%%%%
\subsection{Associative and Commutative Laws}
\label{subsec:logical-associative-commutative}
%%%%%%%%%%%%%%%%%%%%%%%%%%%%%%%%%%%%%%%%%%%%%%%%%%%%%%%%%%%%

Conjunction and disjunction satisfy:

\[
\boxed{
P\land Q
\equiv
Q\land P,
}
\]

\[
\boxed{
P\lor Q
\equiv
Q\lor P.
}
\]

Also,

\[
\boxed{
(P\land Q)\land R
\equiv
P\land(Q\land R),
}
\]

and

\[
\boxed{
(P\lor Q)\lor R
\equiv
P\lor(Q\lor R).
}
\]

%%%%%%%%%%%%%%%%%%%%%%%%%%%%%%%%%%%%%%%%%%%%%%%%%%%%%%%%%%%%
\subsection{Distributive Laws}
\label{subsec:logical-distributive-laws}
%%%%%%%%%%%%%%%%%%%%%%%%%%%%%%%%%%%%%%%%%%%%%%%%%%%%%%%%%%%%

Conjunction distributes over disjunction:

\[
\boxed{
P\land(Q\lor R)
\equiv
(P\land Q)
\lor
(P\land R).
}
\]

Disjunction distributes over conjunction:

\[
\boxed{
P\lor(Q\land R)
\equiv
(P\lor Q)
\land
(P\lor R).
}
\]

%%%%%%%%%%%%%%%%%%%%%%%%%%%%%%%%%%%%%%%%%%%%%%%%%%%%%%%%%%%%
\subsection{Conjunctive Normal Form}
\label{subsec:conjunctive-normal-form}
%%%%%%%%%%%%%%%%%%%%%%%%%%%%%%%%%%%%%%%%%%%%%%%%%%%%%%%%%%%%

A literal is either a propositional variable or its negation.

A clause is a disjunction of literals.

A formula is in conjunctive normal form, abbreviated CNF, if it is a
conjunction of clauses.

For example,

\[
\boxed{
(P\lor\neg Q\lor R)
\land
(\neg P\lor Q).
}
\]

%%%%%%%%%%%%%%%%%%%%%%%%%%%%%%%%%%%%%%%%%%%%%%%%%%%%%%%%%%%%
\subsection{Disjunctive Normal Form}
\label{subsec:disjunctive-normal-form}
%%%%%%%%%%%%%%%%%%%%%%%%%%%%%%%%%%%%%%%%%%%%%%%%%%%%%%%%%%%%

A formula is in disjunctive normal form, abbreviated DNF, if it is a
disjunction of conjunctions of literals.

For example,

\[
\boxed{
(P\land Q)
\lor
(\neg P\land R).
}
\]

Every propositional formula is logically equivalent to a formula in CNF
and to one in DNF.

%%%%%%%%%%%%%%%%%%%%%%%%%%%%%%%%%%%%%%%%%%%%%%%%%%%%%%%%%%%%
\subsection{Worked Example: Conversion to CNF}
\label{subsec:worked-cnf}
%%%%%%%%%%%%%%%%%%%%%%%%%%%%%%%%%%%%%%%%%%%%%%%%%%%%%%%%%%%%

\begin{workedexamplebox}{Eliminating an Implication}

Consider

\[
P\rightarrow(Q\land R).
\]

Eliminate the implication:

\[
\neg P\lor(Q\land R).
\]

Distribute:

\[
(\neg P\lor Q)
\land
(\neg P\lor R).
\]

\begin{answerbox}
\[
\boxed{
(P\rightarrow(Q\land R))
\equiv
(\neg P\lor Q)
\land
(\neg P\lor R).
}
\]
\end{answerbox}

The right-hand side is in conjunctive normal form.

\end{workedexamplebox}

%%%%%%%%%%%%%%%%%%%%%%%%%%%%%%%%%%%%%%%%%%%%%%%%%%%%%%%%%%%%
\subsection{Boolean Functions}
\label{subsec:boolean-functions-formal-logic}
%%%%%%%%%%%%%%%%%%%%%%%%%%%%%%%%%%%%%%%%%%%%%%%%%%%%%%%%%%%%

A Boolean function is a function

\[
\boxed{
f:
\{0,1\}^n
\rightarrow
\{0,1\}.
}
\]

Every propositional formula involving \(n\) variables defines a Boolean
function.

Conversely, every Boolean function can be represented by a propositional
formula.

%%%%%%%%%%%%%%%%%%%%%%%%%%%%%%%%%%%%%%%%%%%%%%%%%%%%%%%%%%%%
\subsection{Functional Completeness}
\label{subsec:functional-completeness}
%%%%%%%%%%%%%%%%%%%%%%%%%%%%%%%%%%%%%%%%%%%%%%%%%%%%%%%%%%%%

A set of logical connectives is functionally complete if every Boolean
function can be expressed using only those connectives.

For example,

\[
\boxed{
\{\neg,\land\}
}
\]

is functionally complete.

Since

\[
P\lor Q
\equiv
\neg(\neg P\land\neg Q),
\]

disjunction can be built from negation and conjunction.

%%%%%%%%%%%%%%%%%%%%%%%%%%%%%%%%%%%%%%%%%%%%%%%%%%%%%%%%%%%%
\subsection{NAND and NOR}
\label{subsec:nand-nor}
%%%%%%%%%%%%%%%%%%%%%%%%%%%%%%%%%%%%%%%%%%%%%%%%%%%%%%%%%%%%

The NAND connective is

\[
\boxed{
P\uparrow Q
=
\neg(P\land Q).
}
\]

The NOR connective is

\[
\boxed{
P\downarrow Q
=
\neg(P\lor Q).
}
\]

Each connective alone is functionally complete.

Thus every Boolean formula can, in principle, be written using only NAND
or only NOR.

%%%%%%%%%%%%%%%%%%%%%%%%%%%%%%%%%%%%%%%%%%%%%%%%%%%%%%%%%%%%
\subsection{The Satisfiability Problem}
\label{subsec:sat-problem}
%%%%%%%%%%%%%%%%%%%%%%%%%%%%%%%%%%%%%%%%%%%%%%%%%%%%%%%%%%%%

The propositional satisfiability problem asks:

\[
\boxed{
\text{Given formula }\varphi,
\text{ does some truth assignment satisfy }\varphi?
}
\]

This problem is abbreviated SAT.

A CNF version is often called CNF-SAT.

SAT is fundamental in:

\[
\boxed{
\text{logic},
}
\]

\[
\boxed{
\text{computer science},
}
\]

\[
\boxed{
\text{verification},
}
\]

and

\[
\boxed{
\text{combinatorial optimization}.
}
\]

%%%%%%%%%%%%%%%%%%%%%%%%%%%%%%%%%%%%%%%%%%%%%%%%%%%%%%%%%%%%
\subsection{Brute-Force Satisfiability}
\label{subsec:brute-force-sat}
%%%%%%%%%%%%%%%%%%%%%%%%%%%%%%%%%%%%%%%%%%%%%%%%%%%%%%%%%%%%

A formula with \(n\) propositional variables has

\[
\boxed{
2^n
}
\]

possible truth assignments.

A naive algorithm evaluates the formula under all assignments.

This is correct but grows exponentially with \(n\).

%%%%%%%%%%%%%%%%%%%%%%%%%%%%%%%%%%%%%%%%%%%%%%%%%%%%%%%%%%%%
\subsection{Formal Proof Systems}
\label{subsec:formal-proof-systems}
%%%%%%%%%%%%%%%%%%%%%%%%%%%%%%%%%%%%%%%%%%%%%%%%%%%%%%%%%%%%

Semantic truth is not the same as formal proof.

A proof system specifies:

\[
\boxed{
\text{axioms}
}
\]

and

\[
\boxed{
\text{inference rules}.
}
\]

A formal derivation is a finite sequence of formulas generated according
to these rules.

%%%%%%%%%%%%%%%%%%%%%%%%%%%%%%%%%%%%%%%%%%%%%%%%%%%%%%%%%%%%
\subsection{Syntactic Derivability}
\label{subsec:syntactic-derivability}
%%%%%%%%%%%%%%%%%%%%%%%%%%%%%%%%%%%%%%%%%%%%%%%%%%%%%%%%%%%%

Write

\[
\boxed{
\Gamma\vdash\varphi
}
\]

when formula \(\varphi\) can be formally derived from assumptions
\(\Gamma\).

The symbol

\[
\vdash
\]

belongs to syntax.

Compare this with

\[
\boxed{
\Gamma\models\varphi,
}
\]

which belongs to semantics.

%%%%%%%%%%%%%%%%%%%%%%%%%%%%%%%%%%%%%%%%%%%%%%%%%%%%%%%%%%%%
\subsection{The Two Turnstiles}
\label{subsec:two-turnstiles}
%%%%%%%%%%%%%%%%%%%%%%%%%%%%%%%%%%%%%%%%%%%%%%%%%%%%%%%%%%%%

The distinction

\[
\boxed{
\vdash
}
\]

versus

\[
\boxed{
\models
}
\]

is fundamental.

\[
\Gamma\vdash\varphi
\]

means:

\[
\boxed{
\text{\(\varphi\) has a formal proof from \(\Gamma\).}
}
\]

\[
\Gamma\models\varphi
\]

means:

\[
\boxed{
\text{every model of \(\Gamma\) makes \(\varphi\) true.}
}
\]

%%%%%%%%%%%%%%%%%%%%%%%%%%%%%%%%%%%%%%%%%%%%%%%%%%%%%%%%%%%%
\subsection{Inference Rules}
\label{subsec:formal-inference-rules}
%%%%%%%%%%%%%%%%%%%%%%%%%%%%%%%%%%%%%%%%%%%%%%%%%%%%%%%%%%%%

An inference rule specifies how formulas may be derived.

For example, modus ponens is

\[
\boxed{
\frac{
\varphi
\qquad
\varphi\rightarrow\psi
}{
\psi
}.
}
\]

If both formulas above the line have been derived, the formula below the
line may be inferred.

%%%%%%%%%%%%%%%%%%%%%%%%%%%%%%%%%%%%%%%%%%%%%%%%%%%%%%%%%%%%
\subsection{Natural Deduction}
\label{subsec:natural-deduction}
%%%%%%%%%%%%%%%%%%%%%%%%%%%%%%%%%%%%%%%%%%%%%%%%%%%%%%%%%%%%

Natural deduction organizes logical reasoning using introduction and
elimination rules for each connective.

For conjunction introduction,

\[
\boxed{
\frac{
\varphi
\qquad
\psi
}{
\varphi\land\psi
}.
}
\]

For conjunction elimination,

\[
\boxed{
\frac{
\varphi\land\psi
}{
\varphi
},
}
\]

and

\[
\boxed{
\frac{
\varphi\land\psi
}{
\psi
}.
}
\]

%%%%%%%%%%%%%%%%%%%%%%%%%%%%%%%%%%%%%%%%%%%%%%%%%%%%%%%%%%%%
\subsection{Implication Introduction}
\label{subsec:implication-introduction}
%%%%%%%%%%%%%%%%%%%%%%%%%%%%%%%%%%%%%%%%%%%%%%%%%%%%%%%%%%%%

To prove

\[
\varphi\rightarrow\psi,
\]

one may temporarily assume \(\varphi\), derive \(\psi\), and discharge the
assumption.

Schematically,

\[
\boxed{
\frac{
[\varphi]
\quad
\vdots
\quad
\psi
}{
\varphi\rightarrow\psi
}.
}
\]

This formalizes conditional proof.

%%%%%%%%%%%%%%%%%%%%%%%%%%%%%%%%%%%%%%%%%%%%%%%%%%%%%%%%%%%%
\subsection{Negation and Contradiction}
\label{subsec:formal-negation-contradiction}
%%%%%%%%%%%%%%%%%%%%%%%%%%%%%%%%%%%%%%%%%%%%%%%%%%%%%%%%%%%%

Let

\[
\bot
\]

denote contradiction or falsity.

Negation may be viewed as

\[
\boxed{
\neg\varphi
\equiv
\varphi\rightarrow\bot.
}
\]

To establish \(\neg\varphi\), assume \(\varphi\) and derive \(\bot\).

%%%%%%%%%%%%%%%%%%%%%%%%%%%%%%%%%%%%%%%%%%%%%%%%%%%%%%%%%%%%
\subsection{Formal Proof by Contradiction}
\label{subsec:formal-proof-by-contradiction}
%%%%%%%%%%%%%%%%%%%%%%%%%%%%%%%%%%%%%%%%%%%%%%%%%%%%%%%%%%%%

In classical logic, to prove \(\varphi\), one may assume

\[
\neg\varphi
\]

and derive contradiction:

\[
\boxed{
\frac{
[\neg\varphi]
\quad
\vdots
\quad
\bot
}{
\varphi
}.
}
\]

This principle depends on the underlying logical system.

Constructive logics treat some classical principles differently.

%%%%%%%%%%%%%%%%%%%%%%%%%%%%%%%%%%%%%%%%%%%%%%%%%%%%%%%%%%%%
\subsection{Deduction Theorem}
\label{subsec:deduction-theorem-formal}
%%%%%%%%%%%%%%%%%%%%%%%%%%%%%%%%%%%%%%%%%%%%%%%%%%%%%%%%%%%%

In many propositional proof systems, a deduction theorem states
schematically that

\[
\boxed{
\Gamma\cup\{\varphi\}
\vdash
\psi
}
\]

is related to

\[
\boxed{
\Gamma
\vdash
\varphi\rightarrow\psi.
}
\]

The exact formulation depends on the proof system.

Natural deduction incorporates this principle directly through implication
introduction.

%%%%%%%%%%%%%%%%%%%%%%%%%%%%%%%%%%%%%%%%%%%%%%%%%%%%%%%%%%%%
\subsection{Soundness}
\label{subsec:formal-soundness}
%%%%%%%%%%%%%%%%%%%%%%%%%%%%%%%%%%%%%%%%%%%%%%%%%%%%%%%%%%%%

\begin{definitionbox}{Soundness}
A proof system is sound if

\[
\boxed{
\Gamma\vdash\varphi
\Longrightarrow
\Gamma\models\varphi.
}
\]
\end{definitionbox}

Thus anything formally provable from assumptions must be semantically
valid in every model satisfying those assumptions.

A sound proof system does not prove false semantic consequences.

%%%%%%%%%%%%%%%%%%%%%%%%%%%%%%%%%%%%%%%%%%%%%%%%%%%%%%%%%%%%
\subsection{Completeness}
\label{subsec:formal-completeness}
%%%%%%%%%%%%%%%%%%%%%%%%%%%%%%%%%%%%%%%%%%%%%%%%%%%%%%%%%%%%

\begin{definitionbox}{Completeness}
A proof system is complete if

\[
\boxed{
\Gamma\models\varphi
\Longrightarrow
\Gamma\vdash\varphi.
}
\]
\end{definitionbox}

Thus every semantic consequence can be formally proved.

Soundness and completeness together give

\[
\boxed{
\Gamma\vdash\varphi
\iff
\Gamma\models\varphi.
}
\]

for the proof system under consideration.

%%%%%%%%%%%%%%%%%%%%%%%%%%%%%%%%%%%%%%%%%%%%%%%%%%%%%%%%%%%%
\subsection{Consistency}
\label{subsec:formal-consistency}
%%%%%%%%%%%%%%%%%%%%%%%%%%%%%%%%%%%%%%%%%%%%%%%%%%%%%%%%%%%%

A theory is syntactically consistent if it does not derive a contradiction.

One formulation is:

\[
\boxed{
T\nvdash\bot.
}
\]

Equivalently, in standard classical settings, there is no formula
\(\varphi\) such that both

\[
T\vdash\varphi
\]

and

\[
T\vdash\neg\varphi.
\]

%%%%%%%%%%%%%%%%%%%%%%%%%%%%%%%%%%%%%%%%%%%%%%%%%%%%%%%%%%%%
\subsection{Explosion}
\label{subsec:principle-explosion}
%%%%%%%%%%%%%%%%%%%%%%%%%%%%%%%%%%%%%%%%%%%%%%%%%%%%%%%%%%%%

In classical logic,

\[
\boxed{
\bot
\vdash
\varphi
}
\]

for arbitrary \(\varphi\).

This is the principle of explosion.

Thus if a classical formal theory is inconsistent, every formula becomes
derivable.

This explains why consistency is foundationally important.

%%%%%%%%%%%%%%%%%%%%%%%%%%%%%%%%%%%%%%%%%%%%%%%%%%%%%%%%%%%%
\subsection{First-Order Logic}
\label{subsec:first-order-logic-introduction}
%%%%%%%%%%%%%%%%%%%%%%%%%%%%%%%%%%%%%%%%%%%%%%%%%%%%%%%%%%%%

Propositional logic treats entire statements as indivisible units.

First-order logic introduces internal mathematical structure.

It can express statements about:

\[
\boxed{
\text{objects},
}
\]

\[
\boxed{
\text{properties},
}
\]

\[
\boxed{
\text{relations},
}
\]

and

\[
\boxed{
\text{functions}.
}
\]

%%%%%%%%%%%%%%%%%%%%%%%%%%%%%%%%%%%%%%%%%%%%%%%%%%%%%%%%%%%%
\subsection{First-Order Signature}
\label{subsec:first-order-signature}
%%%%%%%%%%%%%%%%%%%%%%%%%%%%%%%%%%%%%%%%%%%%%%%%%%%%%%%%%%%%

A first-order signature or vocabulary may contain:

\[
\boxed{
\text{constant symbols},
}
\]

\[
\boxed{
\text{function symbols},
}
\]

and

\[
\boxed{
\text{relation symbols}.
}
\]

Each function and relation symbol has a specified arity.

For example, a group language may contain:

\[
\boxed{
e,
\qquad
\cdot,
\qquad
^{-1}.
}
\]

%%%%%%%%%%%%%%%%%%%%%%%%%%%%%%%%%%%%%%%%%%%%%%%%%%%%%%%%%%%%
\subsection{Variables and Constants}
\label{subsec:first-order-variables-constants}
%%%%%%%%%%%%%%%%%%%%%%%%%%%%%%%%%%%%%%%%%%%%%%%%%%%%%%%%%%%%

Variables are typically written

\[
\boxed{
x,y,z,x_1,x_2,\ldots
}
\]

and range over objects in the domain of a structure.

Constant symbols name distinguished elements.

For example, in arithmetic:

\[
\boxed{
0
}
\]

can be a constant symbol.

%%%%%%%%%%%%%%%%%%%%%%%%%%%%%%%%%%%%%%%%%%%%%%%%%%%%%%%%%%%%
\subsection{Function Symbols}
\label{subsec:first-order-function-symbols}
%%%%%%%%%%%%%%%%%%%%%%%%%%%%%%%%%%%%%%%%%%%%%%%%%%%%%%%%%%%%

An \(n\)-ary function symbol represents an operation taking \(n\)
arguments.

For example,

\[
\boxed{
+
}
\]

may be a binary function symbol.

A unary function symbol may represent negation:

\[
\boxed{
-.
}
\]

These are syntactic symbols until interpreted inside a structure.

%%%%%%%%%%%%%%%%%%%%%%%%%%%%%%%%%%%%%%%%%%%%%%%%%%%%%%%%%%%%
\subsection{Relation Symbols}
\label{subsec:first-order-relation-symbols}
%%%%%%%%%%%%%%%%%%%%%%%%%%%%%%%%%%%%%%%%%%%%%%%%%%%%%%%%%%%%

An \(n\)-ary relation symbol expresses a property relating \(n\) objects.

Examples include:

\[
\boxed{
<
}
\]

as a binary relation, and

\[
\boxed{
=
}
\]

as equality.

A unary relation symbol can represent a property such as

\[
\boxed{
P(x).
}
\]

%%%%%%%%%%%%%%%%%%%%%%%%%%%%%%%%%%%%%%%%%%%%%%%%%%%%%%%%%%%%
\subsection{Terms}
\label{subsec:first-order-terms}
%%%%%%%%%%%%%%%%%%%%%%%%%%%%%%%%%%%%%%%%%%%%%%%%%%%%%%%%%%%%

Terms represent objects.

They are defined recursively:

\begin{enumerate}
    \item every variable is a term;
    \item every constant symbol is a term;
    \item if \(f\) is an \(n\)-ary function symbol and
    \[
    t_1,\ldots,t_n
    \]
    are terms, then
    \[
    \boxed{
    f(t_1,\ldots,t_n)
    }
    \]
    is a term.
\end{enumerate}

%%%%%%%%%%%%%%%%%%%%%%%%%%%%%%%%%%%%%%%%%%%%%%%%%%%%%%%%%%%%
\subsection{Examples of Terms}
\label{subsec:first-order-term-examples}
%%%%%%%%%%%%%%%%%%%%%%%%%%%%%%%%%%%%%%%%%%%%%%%%%%%%%%%%%%%%

In the language of arithmetic,

\[
x,
\]

\[
0,
\]

\[
x+1,
\]

and

\[
(x+y)\cdot z
\]

are terms.

A term does not by itself assert a proposition.

It denotes an object after interpretation.

%%%%%%%%%%%%%%%%%%%%%%%%%%%%%%%%%%%%%%%%%%%%%%%%%%%%%%%%%%%%
\subsection{Atomic Formulas}
\label{subsec:atomic-formulas}
%%%%%%%%%%%%%%%%%%%%%%%%%%%%%%%%%%%%%%%%%%%%%%%%%%%%%%%%%%%%

If

\[
R
\]

is an \(n\)-ary relation symbol and

\[
t_1,\ldots,t_n
\]

are terms, then

\[
\boxed{
R(t_1,\ldots,t_n)
}
\]

is an atomic formula.

Equality produces atomic formulas of the form

\[
\boxed{
t_1=t_2.
}
\]

%%%%%%%%%%%%%%%%%%%%%%%%%%%%%%%%%%%%%%%%%%%%%%%%%%%%%%%%%%%%
\subsection{First-Order Formulas}
\label{subsec:first-order-formulas}
%%%%%%%%%%%%%%%%%%%%%%%%%%%%%%%%%%%%%%%%%%%%%%%%%%%%%%%%%%%%

Starting with atomic formulas, construct more formulas using

\[
\boxed{
\neg,
\land,
\lor,
\rightarrow,
\leftrightarrow,
}
\]

and the quantifiers

\[
\boxed{
\forall,
\qquad
\exists.
}
\]

Thus if \(\varphi\) is a formula and \(x\) is a variable, then

\[
\boxed{
\forall x\,\varphi
}
\]

and

\[
\boxed{
\exists x\,\varphi
}
\]

are formulas.

%%%%%%%%%%%%%%%%%%%%%%%%%%%%%%%%%%%%%%%%%%%%%%%%%%%%%%%%%%%%
\subsection{Free and Bound Variables}
\label{subsec:free-bound-variables}
%%%%%%%%%%%%%%%%%%%%%%%%%%%%%%%%%%%%%%%%%%%%%%%%%%%%%%%%%%%%

In

\[
\forall x\,P(x,y),
\]

the variable \(x\) is bound by the quantifier.

The variable \(y\) remains free.

Thus the formula depends on a value assigned to \(y\).

A variable occurrence is free when it is not within the scope of a
quantifier binding that variable.

%%%%%%%%%%%%%%%%%%%%%%%%%%%%%%%%%%%%%%%%%%%%%%%%%%%%%%%%%%%%
\subsection{Sentences}
\label{subsec:first-order-sentences}
%%%%%%%%%%%%%%%%%%%%%%%%%%%%%%%%%%%%%%%%%%%%%%%%%%%%%%%%%%%%

\begin{definitionbox}{Sentence}
A sentence is a formula with no free variables.
\end{definitionbox}

For example,

\[
\boxed{
\forall x
\,
(x=x)
}
\]

is a sentence.

By contrast,

\[
\boxed{
x=x
}
\]

contains a free variable.

%%%%%%%%%%%%%%%%%%%%%%%%%%%%%%%%%%%%%%%%%%%%%%%%%%%%%%%%%%%%
\subsection{Scope of a Quantifier}
\label{subsec:quantifier-scope}
%%%%%%%%%%%%%%%%%%%%%%%%%%%%%%%%%%%%%%%%%%%%%%%%%%%%%%%%%%%%

In

\[
\forall x
\,
(P(x)\rightarrow Q(x)),
\]

the scope of

\[
\forall x
\]

is the formula

\[
P(x)\rightarrow Q(x).
\]

Precise scope is essential when formulas contain several quantifiers.

For example,

\[
\forall x\exists y\,R(x,y)
\]

and

\[
\exists y\forall x\,R(x,y)
\]

generally have different meanings.

%%%%%%%%%%%%%%%%%%%%%%%%%%%%%%%%%%%%%%%%%%%%%%%%%%%%%%%%%%%%
\subsection{Worked Example: Quantifier Order}
\label{subsec:worked-quantifier-order}
%%%%%%%%%%%%%%%%%%%%%%%%%%%%%%%%%%%%%%%%%%%%%%%%%%%%%%%%%%%%

\begin{workedexamplebox}{Different Quantifier Orders}

Over the real numbers, consider

\[
\forall x\exists y
\,
(y>x).
\]

This is true because for any real \(x\), one may choose

\[
y=x+1.
\]

Now consider

\[
\exists y\forall x
\,
(y>x).
\]

This would require one real number larger than every real number, which
does not exist.

\begin{answerbox}
\[
\boxed{
\forall x\exists y\,(y>x)
\text{ is true,}
}
\]

while

\[
\boxed{
\exists y\forall x\,(y>x)
\text{ is false over }\mathbb{R}.
}
\]
\end{answerbox}

\end{workedexamplebox}

%%%%%%%%%%%%%%%%%%%%%%%%%%%%%%%%%%%%%%%%%%%%%%%%%%%%%%%%%%%%
\subsection{Substitution}
\label{subsec:first-order-substitution}
%%%%%%%%%%%%%%%%%%%%%%%%%%%%%%%%%%%%%%%%%%%%%%%%%%%%%%%%%%%%

Substitution replaces a free occurrence of variable \(x\) by a term \(t\).

Write schematically

\[
\boxed{
\varphi[t/x].
}
\]

Substitution must respect quantifier scope.

A naive textual replacement can change the meaning of a formula.

%%%%%%%%%%%%%%%%%%%%%%%%%%%%%%%%%%%%%%%%%%%%%%%%%%%%%%%%%%%%
\subsection{Variable Capture}
\label{subsec:variable-capture}
%%%%%%%%%%%%%%%%%%%%%%%%%%%%%%%%%%%%%%%%%%%%%%%%%%%%%%%%%%%%

Consider

\[
\exists y\,R(x,y).
\]

If one substitutes term \(y\) for free variable \(x\) naively, the result
would be

\[
\exists y\,R(y,y).
\]

The substituted \(y\) has become bound by the existing quantifier.

This is called variable capture.

Formal substitution rules prevent this unintended change.

%%%%%%%%%%%%%%%%%%%%%%%%%%%%%%%%%%%%%%%%%%%%%%%%%%%%%%%%%%%%
\subsection{Renaming Bound Variables}
\label{subsec:renaming-bound-variables}
%%%%%%%%%%%%%%%%%%%%%%%%%%%%%%%%%%%%%%%%%%%%%%%%%%%%%%%%%%%%

Bound variables can often be renamed without changing meaning.

For example,

\[
\boxed{
\forall x\,P(x)
}
\]

and

\[
\boxed{
\forall z\,P(z)
}
\]

are equivalent when the renaming avoids variable capture.

This is sometimes called alpha-renaming.

%%%%%%%%%%%%%%%%%%%%%%%%%%%%%%%%%%%%%%%%%%%%%%%%%%%%%%%%%%%%
\subsection{Structures}
\label{subsec:first-order-structures}
%%%%%%%%%%%%%%%%%%%%%%%%%%%%%%%%%%%%%%%%%%%%%%%%%%%%%%%%%%%%

A first-order structure provides meaning to a formal language.

A structure \(\mathcal{M}\) contains:

\[
\boxed{
\text{a nonempty domain }M,
}
\]

together with interpretations of the constant, function, and relation
symbols.

For example,

\[
\boxed{
\mathcal{M}
=
(\mathbb{R},+,\,\cdot,\,<,\,0,\,1)
}
\]

interprets arithmetic symbols over the real numbers.

%%%%%%%%%%%%%%%%%%%%%%%%%%%%%%%%%%%%%%%%%%%%%%%%%%%%%%%%%%%%
\subsection{Interpretation of Constants}
\label{subsec:interpretation-constants}
%%%%%%%%%%%%%%%%%%%%%%%%%%%%%%%%%%%%%%%%%%%%%%%%%%%%%%%%%%%%

If \(c\) is a constant symbol, a structure assigns an element

\[
\boxed{
c^{\mathcal{M}}
\in M.
}
\]

For example, in arithmetic,

\[
0^{\mathcal{M}}
\]

may denote the actual number \(0\).

%%%%%%%%%%%%%%%%%%%%%%%%%%%%%%%%%%%%%%%%%%%%%%%%%%%%%%%%%%%%
\subsection{Interpretation of Functions}
\label{subsec:interpretation-functions}
%%%%%%%%%%%%%%%%%%%%%%%%%%%%%%%%%%%%%%%%%%%%%%%%%%%%%%%%%%%%

If \(f\) is an \(n\)-ary function symbol, then

\[
\boxed{
f^{\mathcal{M}}
:
M^n
\rightarrow
M.
}
\]

For example, the symbol \(+\) may be interpreted by ordinary addition on
\(\mathbb{R}\).

%%%%%%%%%%%%%%%%%%%%%%%%%%%%%%%%%%%%%%%%%%%%%%%%%%%%%%%%%%%%
\subsection{Interpretation of Relations}
\label{subsec:interpretation-relations}
%%%%%%%%%%%%%%%%%%%%%%%%%%%%%%%%%%%%%%%%%%%%%%%%%%%%%%%%%%%%

If \(R\) is an \(n\)-ary relation symbol, then

\[
\boxed{
R^{\mathcal{M}}
\subseteq
M^n.
}
\]

For example,

\[
<^{\mathcal{M}}
\]

may be the usual strict-order relation on \(\mathbb{R}\).

%%%%%%%%%%%%%%%%%%%%%%%%%%%%%%%%%%%%%%%%%%%%%%%%%%%%%%%%%%%%
\subsection{Variable Assignments}
\label{subsec:first-order-variable-assignments}
%%%%%%%%%%%%%%%%%%%%%%%%%%%%%%%%%%%%%%%%%%%%%%%%%%%%%%%%%%%%

A variable assignment is a function

\[
\boxed{
s:
\operatorname{Var}
\rightarrow
M.
}
\]

It assigns an element of the domain to every variable.

Terms containing variables are evaluated relative to both:

\[
\boxed{
\mathcal{M}
}
\]

and

\[
\boxed{
s.
}
\]

%%%%%%%%%%%%%%%%%%%%%%%%%%%%%%%%%%%%%%%%%%%%%%%%%%%%%%%%%%%%
\subsection{Evaluation of Terms}
\label{subsec:evaluation-of-terms}
%%%%%%%%%%%%%%%%%%%%%%%%%%%%%%%%%%%%%%%%%%%%%%%%%%%%%%%%%%%%

The value of term \(t\) in structure \(\mathcal{M}\) under assignment \(s\)
is written conceptually as

\[
\boxed{
t^{\mathcal{M}}[s].
}
\]

Variables are evaluated using \(s\).

Constants are evaluated using the structure.

Compound terms are evaluated recursively through the interpreted function
symbols.

%%%%%%%%%%%%%%%%%%%%%%%%%%%%%%%%%%%%%%%%%%%%%%%%%%%%%%%%%%%%
\subsection{First-Order Satisfaction}
\label{subsec:first-order-satisfaction}
%%%%%%%%%%%%%%%%%%%%%%%%%%%%%%%%%%%%%%%%%%%%%%%%%%%%%%%%%%%%

Write

\[
\boxed{
\mathcal{M},s
\models
\varphi
}
\]

when formula \(\varphi\) is true in structure \(\mathcal{M}\) under
assignment \(s\).

For sentences, the assignment does not affect truth, so one writes

\[
\boxed{
\mathcal{M}
\models
\varphi.
}
\]

%%%%%%%%%%%%%%%%%%%%%%%%%%%%%%%%%%%%%%%%%%%%%%%%%%%%%%%%%%%%
\subsection{Atomic Satisfaction}
\label{subsec:atomic-satisfaction}
%%%%%%%%%%%%%%%%%%%%%%%%%%%%%%%%%%%%%%%%%%%%%%%%%%%%%%%%%%%%

For atomic formula

\[
R(t_1,\ldots,t_n),
\]

we have

\[
\boxed{
\mathcal{M},s
\models
R(t_1,\ldots,t_n)
}
\]

exactly when

\[
\boxed{
(
t_1^{\mathcal{M}}[s],
\ldots,
t_n^{\mathcal{M}}[s]
)
\in
R^{\mathcal{M}}.
}
\]

Thus semantic truth is determined by the interpretation of the relation.

%%%%%%%%%%%%%%%%%%%%%%%%%%%%%%%%%%%%%%%%%%%%%%%%%%%%%%%%%%%%
\subsection{Universal Quantifier Semantics}
\label{subsec:universal-quantifier-semantics}
%%%%%%%%%%%%%%%%%%%%%%%%%%%%%%%%%%%%%%%%%%%%%%%%%%%%%%%%%%%%

The formula

\[
\forall x\,\varphi
\]

is true in \(\mathcal{M}\) under assignment \(s\) when \(\varphi\) is true
for every possible value of \(x\).

Formally,

\[
\boxed{
\mathcal{M},s
\models
\forall x\,\varphi
}
\]

if and only if

\[
\boxed{
\mathcal{M},s[x\mapsto a]
\models
\varphi
}
\]

for every

\[
a\in M.
\]

%%%%%%%%%%%%%%%%%%%%%%%%%%%%%%%%%%%%%%%%%%%%%%%%%%%%%%%%%%%%
\subsection{Existential Quantifier Semantics}
\label{subsec:existential-quantifier-semantics}
%%%%%%%%%%%%%%%%%%%%%%%%%%%%%%%%%%%%%%%%%%%%%%%%%%%%%%%%%%%%

The formula

\[
\exists x\,\varphi
\]

is true if at least one element of the domain makes \(\varphi\) true.

Thus

\[
\boxed{
\mathcal{M},s
\models
\exists x\,\varphi
}
\]

if and only if there exists

\[
a\in M
\]

such that

\[
\boxed{
\mathcal{M},s[x\mapsto a]
\models
\varphi.
}
\]

%%%%%%%%%%%%%%%%%%%%%%%%%%%%%%%%%%%%%%%%%%%%%%%%%%%%%%%%%%%%
\subsection{Negating Quantifiers}
\label{subsec:formal-negating-quantifiers}
%%%%%%%%%%%%%%%%%%%%%%%%%%%%%%%%%%%%%%%%%%%%%%%%%%%%%%%%%%%%

In classical first-order logic,

\[
\boxed{
\neg\forall x\,\varphi
\equiv
\exists x\,\neg\varphi,
}
\]

and

\[
\boxed{
\neg\exists x\,\varphi
\equiv
\forall x\,\neg\varphi.
}
\]

These are quantified analogues of De Morgan's laws.

%%%%%%%%%%%%%%%%%%%%%%%%%%%%%%%%%%%%%%%%%%%%%%%%%%%%%%%%%%%%
\subsection{Worked Example: Negating a Quantified Statement}
\label{subsec:worked-negating-quantifier}
%%%%%%%%%%%%%%%%%%%%%%%%%%%%%%%%%%%%%%%%%%%%%%%%%%%%%%%%%%%%

\begin{workedexamplebox}{Negating a Universal Statement}

Consider

\[
\forall x
\,
P(x).
\]

Its negation is

\[
\neg
\forall x
\,
P(x).
\]

Move the negation through the universal quantifier:

\[
\begin{answerbox}
\[
\boxed{
\neg
\forall x\,P(x)
\equiv
\exists x\,\neg P(x).
}
\]
\end{answerbox}

Thus ``not every \(x\) has property \(P\)'' means ``some \(x\) fails to
have property \(P\).''

\end{workedexamplebox}

%%%%%%%%%%%%%%%%%%%%%%%%%%%%%%%%%%%%%%%%%%%%%%%%%%%%%%%%%%%%
\subsection{Models}
\label{subsec:formal-models}
%%%%%%%%%%%%%%%%%%%%%%%%%%%%%%%%%%%%%%%%%%%%%%%%%%%%%%%%%%%%

A structure \(\mathcal{M}\) is a model of a sentence \(\varphi\) if

\[
\boxed{
\mathcal{M}\models\varphi.
}
\]

For a theory \(T\),

\[
\boxed{
\mathcal{M}\models T
}
\]

means \(\mathcal{M}\) satisfies every sentence in \(T\).

Thus a theory describes a class of mathematical structures.

%%%%%%%%%%%%%%%%%%%%%%%%%%%%%%%%%%%%%%%%%%%%%%%%%%%%%%%%%%%%
\subsection{Theories}
\label{subsec:formal-theories}
%%%%%%%%%%%%%%%%%%%%%%%%%%%%%%%%%%%%%%%%%%%%%%%%%%%%%%%%%%%%

A theory is a set of sentences in a fixed language.

Examples include axiomatizations of:

\[
\boxed{
\text{groups},
}
\]

\[
\boxed{
\text{orders},
}
\]

\[
\boxed{
\text{fields},
}
\]

and

\[
\boxed{
\text{arithmetic}.
}
\]

A theory may be finite or infinite.

%%%%%%%%%%%%%%%%%%%%%%%%%%%%%%%%%%%%%%%%%%%%%%%%%%%%%%%%%%%%
\subsection{Example: Group Axioms in First-Order Form}
\label{subsec:first-order-group-axioms}
%%%%%%%%%%%%%%%%%%%%%%%%%%%%%%%%%%%%%%%%%%%%%%%%%%%%%%%%%%%%

In a language containing binary function symbol \(\cdot\), unary function
symbol \(^{-1}\), and constant \(e\), group axioms can be expressed as
first-order sentences.

Associativity:

\[
\boxed{
\forall x\forall y\forall z
\,
((x\cdot y)\cdot z
=
x\cdot(y\cdot z)).
}
\]

Identity:

\[
\boxed{
\forall x
\,
(e\cdot x=x
\land
x\cdot e=x).
}
\]

Inverse:

\[
\boxed{
\forall x
\,
(x\cdot x^{-1}=e
\land
x^{-1}\cdot x=e).
}
\]

%%%%%%%%%%%%%%%%%%%%%%%%%%%%%%%%%%%%%%%%%%%%%%%%%%%%%%%%%%%%
\subsection{First-Order Validity}
\label{subsec:first-order-validity}
%%%%%%%%%%%%%%%%%%%%%%%%%%%%%%%%%%%%%%%%%%%%%%%%%%%%%%%%%%%%

A sentence \(\varphi\) is logically valid if

\[
\boxed{
\mathcal{M}\models\varphi
}
\]

for every structure \(\mathcal{M}\) for the language.

Write

\[
\boxed{
\models\varphi.
}
\]

This is much stronger than being true in one particular mathematical
structure.

%%%%%%%%%%%%%%%%%%%%%%%%%%%%%%%%%%%%%%%%%%%%%%%%%%%%%%%%%%%%
\subsection{First-Order Satisfiability}
\label{subsec:first-order-satisfiability}
%%%%%%%%%%%%%%%%%%%%%%%%%%%%%%%%%%%%%%%%%%%%%%%%%%%%%%%%%%%%

A sentence \(\varphi\) is satisfiable if there exists at least one structure

\[
\mathcal{M}
\]

such that

\[
\boxed{
\mathcal{M}\models\varphi.
}
\]

A theory \(T\) is satisfiable if it has at least one model.

%%%%%%%%%%%%%%%%%%%%%%%%%%%%%%%%%%%%%%%%%%%%%%%%%%%%%%%%%%%%
\subsection{Semantic Consequence for Theories}
\label{subsec:semantic-consequence-theories}
%%%%%%%%%%%%%%%%%%%%%%%%%%%%%%%%%%%%%%%%%%%%%%%%%%%%%%%%%%%%

For theory \(T\),

\[
\boxed{
T\models\varphi
}
\]

means every model of \(T\) is also a model of \(\varphi\).

Thus \(\varphi\) is true in every structure satisfying the axioms of the
theory.

%%%%%%%%%%%%%%%%%%%%%%%%%%%%%%%%%%%%%%%%%%%%%%%%%%%%%%%%%%%%
\subsection{Example: Group-Theoretic Consequence}
\label{subsec:group-theoretic-consequence}
%%%%%%%%%%%%%%%%%%%%%%%%%%%%%%%%%%%%%%%%%%%%%%%%%%%%%%%%%%%%

Let \(T_{\mathrm{grp}}\) denote the first-order group axioms.

The identity element is unique in every group.

Thus the sentence expressing uniqueness of the identity is a semantic
consequence:

\[
\boxed{
T_{\mathrm{grp}}
\models
\text{``the identity is unique.''}
}
\]

A formal proof system should also be able to derive that consequence.

%%%%%%%%%%%%%%%%%%%%%%%%%%%%%%%%%%%%%%%%%%%%%%%%%%%%%%%%%%%%
\subsection{Equality}
\label{subsec:first-order-equality}
%%%%%%%%%%%%%%%%%%%%%%%%%%%%%%%%%%%%%%%%%%%%%%%%%%%%%%%%%%%%

First-order logic with equality treats

\[
=
\]

as a distinguished binary relation interpreted as actual equality on the
domain.

Basic logical properties include:

\[
\boxed{
\forall x\,(x=x),
}
\]

and substitutivity:

if

\[
x=y,
\]

then formulas and functions should respect replacement of equals by equals
in appropriate contexts.

%%%%%%%%%%%%%%%%%%%%%%%%%%%%%%%%%%%%%%%%%%%%%%%%%%%%%%%%%%%%
\subsection{Equality and Functions}
\label{subsec:equality-functions}
%%%%%%%%%%%%%%%%%%%%%%%%%%%%%%%%%%%%%%%%%%%%%%%%%%%%%%%%%%%%

If

\[
x_1=y_1,
\ldots,
x_n=y_n,
\]

then for an \(n\)-ary function symbol \(f\),

\[
\boxed{
f(x_1,\ldots,x_n)
=
f(y_1,\ldots,y_n).
}
\]

Functions therefore preserve equality.

%%%%%%%%%%%%%%%%%%%%%%%%%%%%%%%%%%%%%%%%%%%%%%%%%%%%%%%%%%%%
\subsection{Equality and Relations}
\label{subsec:equality-relations}
%%%%%%%%%%%%%%%%%%%%%%%%%%%%%%%%%%%%%%%%%%%%%%%%%%%%%%%%%%%%

For relation \(R\), if corresponding arguments are equal, relation truth is
preserved.

Informally,

\[
x_i=y_i
\]

for all \(i\) implies

\[
\boxed{
R(x_1,\ldots,x_n)
\leftrightarrow
R(y_1,\ldots,y_n).
}
\]

This formalizes substitution of equals.

%%%%%%%%%%%%%%%%%%%%%%%%%%%%%%%%%%%%%%%%%%%%%%%%%%%%%%%%%%%%
\subsection{First-Order Proof Systems}
\label{subsec:first-order-proof-systems}
%%%%%%%%%%%%%%%%%%%%%%%%%%%%%%%%%%%%%%%%%%%%%%%%%%%%%%%%%%%%

First-order proof systems extend propositional inference rules with rules
for quantifiers.

Natural-deduction systems include:

\[
\boxed{
\forall\text{-introduction},
}
\]

\[
\boxed{
\forall\text{-elimination},
}
\]

\[
\boxed{
\exists\text{-introduction},
}
\]

and

\[
\boxed{
\exists\text{-elimination}.
}
\]

Side conditions prevent invalid reasoning with free variables.

%%%%%%%%%%%%%%%%%%%%%%%%%%%%%%%%%%%%%%%%%%%%%%%%%%%%%%%%%%%%
\subsection{Universal Elimination}
\label{subsec:universal-elimination}
%%%%%%%%%%%%%%%%%%%%%%%%%%%%%%%%%%%%%%%%%%%%%%%%%%%%%%%%%%%%

From

\[
\forall x\,\varphi(x),
\]

one may infer

\[
\boxed{
\varphi(t)
}
\]

for an appropriate term \(t\).

Schematically,

\[
\boxed{
\frac{
\forall x\,\varphi(x)
}{
\varphi(t)
}.
}
\]

This formalizes specialization.

%%%%%%%%%%%%%%%%%%%%%%%%%%%%%%%%%%%%%%%%%%%%%%%%%%%%%%%%%%%%
\subsection{Universal Introduction}
\label{subsec:universal-introduction}
%%%%%%%%%%%%%%%%%%%%%%%%%%%%%%%%%%%%%%%%%%%%%%%%%%%%%%%%%%%%

To conclude

\[
\forall x\,\varphi(x),
\]

one must establish \(\varphi(x)\) for an arbitrary \(x\).

The variable must not depend on special assumptions that restrict its
value.

This side condition prevents invalid generalization.

%%%%%%%%%%%%%%%%%%%%%%%%%%%%%%%%%%%%%%%%%%%%%%%%%%%%%%%%%%%%
\subsection{Existential Introduction}
\label{subsec:existential-introduction}
%%%%%%%%%%%%%%%%%%%%%%%%%%%%%%%%%%%%%%%%%%%%%%%%%%%%%%%%%%%%

From

\[
\varphi(t),
\]

one may infer

\[
\boxed{
\exists x\,\varphi(x).
}
\]

This formalizes the idea that exhibiting one witness proves existence.

%%%%%%%%%%%%%%%%%%%%%%%%%%%%%%%%%%%%%%%%%%%%%%%%%%%%%%%%%%%%
\subsection{Existential Elimination}
\label{subsec:existential-elimination}
%%%%%%%%%%%%%%%%%%%%%%%%%%%%%%%%%%%%%%%%%%%%%%%%%%%%%%%%%%%%

From

\[
\exists x\,\varphi(x),
\]

one may reason using a fresh representative satisfying \(\varphi\).

However, the final conclusion must not depend on any special property of
that representative beyond \(\varphi\).

The freshness condition is essential.

%%%%%%%%%%%%%%%%%%%%%%%%%%%%%%%%%%%%%%%%%%%%%%%%%%%%%%%%%%%%
\subsection{Soundness of First-Order Logic}
\label{subsec:first-order-soundness}
%%%%%%%%%%%%%%%%%%%%%%%%%%%%%%%%%%%%%%%%%%%%%%%%%%%%%%%%%%%%

A sound first-order proof system satisfies

\[
\boxed{
T\vdash\varphi
\Longrightarrow
T\models\varphi.
}
\]

Proof rules preserve truth in every structure.

Soundness therefore connects syntactic deduction to semantic validity.

%%%%%%%%%%%%%%%%%%%%%%%%%%%%%%%%%%%%%%%%%%%%%%%%%%%%%%%%%%%%
\subsection{Completeness of First-Order Logic}
\label{subsec:first-order-completeness}
%%%%%%%%%%%%%%%%%%%%%%%%%%%%%%%%%%%%%%%%%%%%%%%%%%%%%%%%%%%%

Gödel's completeness theorem for first-order logic states, in one standard
form, that

\[
\boxed{
T\models\varphi
\Longrightarrow
T\vdash\varphi.
}
\]

Together with soundness,

\[
\boxed{
T\vdash\varphi
\iff
T\models\varphi.
}
\]

This concerns completeness of the logical proof system.

It should not be confused with Gödel's incompleteness theorems for certain
formal theories, which concern a different notion.

%%%%%%%%%%%%%%%%%%%%%%%%%%%%%%%%%%%%%%%%%%%%%%%%%%%%%%%%%%%%
\subsection{Syntactic Consistency and Models}
\label{subsec:consistency-and-models}
%%%%%%%%%%%%%%%%%%%%%%%%%%%%%%%%%%%%%%%%%%%%%%%%%%%%%%%%%%%%

A standard consequence of first-order soundness and completeness is that
a theory is syntactically consistent exactly when it has a model, under the
usual definitions and proof systems.

Schematically,

\[
\boxed{
T
\text{ consistent}
\iff
T
\text{ satisfiable}.
}
\]

This establishes a deep connection between proof and structure.

%%%%%%%%%%%%%%%%%%%%%%%%%%%%%%%%%%%%%%%%%%%%%%%%%%%%%%%%%%%%
\subsection{Compactness Theorem Preview}
\label{subsec:compactness-preview}
%%%%%%%%%%%%%%%%%%%%%%%%%%%%%%%%%%%%%%%%%%%%%%%%%%%%%%%%%%%%

The first-order compactness theorem states:

\[
\boxed{
T
\text{ has a model}
}
\]

if and only if every finite subset of \(T\) has a model.

Equivalently, if

\[
T\models\varphi,
\]

then some finite subset

\[
T_0\subseteq T
\]

already satisfies

\[
\boxed{
T_0\models\varphi.
}
\]

Compactness has many surprising consequences.

%%%%%%%%%%%%%%%%%%%%%%%%%%%%%%%%%%%%%%%%%%%%%%%%%%%%%%%%%%%%
\subsection{Worked Example: Compactness Intuition}
\label{subsec:worked-compactness-intuition}
%%%%%%%%%%%%%%%%%%%%%%%%%%%%%%%%%%%%%%%%%%%%%%%%%%%%%%%%%%%%

\begin{workedexamplebox}{Forcing an Infinite Model}

Suppose a first-order theory contains sentences asserting:

\[
\text{there are at least }1\text{ elements},
\]

\[
\text{there are at least }2\text{ elements},
\]

\[
\text{there are at least }3\text{ elements},
\]

and so on.

Every finite subset requires only some sufficiently large finite model.

Therefore every finite subset is satisfiable.

By compactness, the entire theory has a model.

But any model satisfying every sentence must contain infinitely many
elements.

\begin{answerbox}
\[
\boxed{
\text{Compactness guarantees an infinite model.}
}
\]
\end{answerbox}

\end{workedexamplebox}

%%%%%%%%%%%%%%%%%%%%%%%%%%%%%%%%%%%%%%%%%%%%%%%%%%%%%%%%%%%%
\subsection{Löwenheim--Skolem Phenomena Preview}
\label{subsec:lowenheim-skolem-preview}
%%%%%%%%%%%%%%%%%%%%%%%%%%%%%%%%%%%%%%%%%%%%%%%%%%%%%%%%%%%%

First-order logic has strong restrictions on which infinite cardinalities
it can distinguish.

Roughly, if a suitable first-order theory has an infinite model, then under
standard hypotheses it often has models of other infinite cardinalities as
well.

The downward and upward Löwenheim--Skolem theorems make this precise.

These results are developed more fully in model theory.

%%%%%%%%%%%%%%%%%%%%%%%%%%%%%%%%%%%%%%%%%%%%%%%%%%%%%%%%%%%%
\subsection{Expressive Power}
\label{subsec:first-order-expressive-power}
%%%%%%%%%%%%%%%%%%%%%%%%%%%%%%%%%%%%%%%%%%%%%%%%%%%%%%%%%%%%

First-order logic can express many mathematical properties.

For example:

\[
\boxed{
\forall x\forall y
\,
(x+y=y+x)
}
\]

expresses commutativity.

It can express:

\[
\boxed{
\text{group axioms},
}
\]

\[
\boxed{
\text{field axioms},
}
\]

and many graph-theoretic properties.

However, not every mathematical property is first-order expressible.

%%%%%%%%%%%%%%%%%%%%%%%%%%%%%%%%%%%%%%%%%%%%%%%%%%%%%%%%%%%%
\subsection{Limitations of First-Order Expressibility}
\label{subsec:first-order-limitations}
%%%%%%%%%%%%%%%%%%%%%%%%%%%%%%%%%%%%%%%%%%%%%%%%%%%%%%%%%%%%

Within ordinary first-order logic, one cannot simply quantify over arbitrary
subsets, functions, or relations unless they are already elements of the
domain.

First-order quantifiers range over individual domain elements:

\[
\boxed{
\forall x,
\qquad
\exists x.
}
\]

Second-order logic extends quantification to relations or subsets.

This increases expressive power but changes important metatheoretic
properties.

%%%%%%%%%%%%%%%%%%%%%%%%%%%%%%%%%%%%%%%%%%%%%%%%%%%%%%%%%%%%
\subsection{First-Order Versus Second-Order Statements}
\label{subsec:first-vs-second-order}
%%%%%%%%%%%%%%%%%%%%%%%%%%%%%%%%%%%%%%%%%%%%%%%%%%%%%%%%%%%%

A first-order statement quantifies over elements.

For example,

\[
\boxed{
\forall x\exists y\,(x<y).
}
\]

A second-order language may quantify over subsets or relations.

For example, mathematical induction can be expressed categorically using a
quantifier over all subsets of the natural numbers.

This distinction is important in foundations.

%%%%%%%%%%%%%%%%%%%%%%%%%%%%%%%%%%%%%%%%%%%%%%%%%%%%%%%%%%%%
\subsection{Translation Into Formal Logic}
\label{subsec:translation-formal-logic}
%%%%%%%%%%%%%%%%%%%%%%%%%%%%%%%%%%%%%%%%%%%%%%%%%%%%%%%%%%%%

A major skill is translating ordinary mathematical statements into formal
language.

For example,

\[
\boxed{
\text{``Every integer has an additive inverse.''}
}
\]

can be represented schematically as

\[
\boxed{
\forall x
\exists y
\,
(x+y=0).
}
\]

The intended structure determines the domain of quantification.

%%%%%%%%%%%%%%%%%%%%%%%%%%%%%%%%%%%%%%%%%%%%%%%%%%%%%%%%%%%%
\subsection{Uniqueness}
\label{subsec:formal-uniqueness}
%%%%%%%%%%%%%%%%%%%%%%%%%%%%%%%%%%%%%%%%%%%%%%%%%%%%%%%%%%%%

The statement

\[
\exists!x\,P(x)
\]

means there exists exactly one \(x\) satisfying \(P\).

It can be expanded as

\[
\boxed{
\exists x
\left(
P(x)
\land
\forall y
\left(
P(y)\rightarrow y=x
\right)
\right).
}
\]

Thus uniqueness does not require a new fundamental quantifier.

%%%%%%%%%%%%%%%%%%%%%%%%%%%%%%%%%%%%%%%%%%%%%%%%%%%%%%%%%%%%
\subsection{Worked Example: Formalizing Uniqueness}
\label{subsec:worked-formalizing-uniqueness}
%%%%%%%%%%%%%%%%%%%%%%%%%%%%%%%%%%%%%%%%%%%%%%%%%%%%%%%%%%%%

\begin{workedexamplebox}{Unique Additive Identity}

The statement

\[
\text{``There is exactly one additive identity''}
\]

can be written

\[
\exists e
\left[
\forall x
\,
(e+x=x)
\land
\forall f
\left(
\forall x\,(f+x=x)
\rightarrow
f=e
\right)
\right].
\]

\begin{answerbox}
\[
\boxed{
\exists e
\left[
\forall x\,(e+x=x)
\land
\forall f
\left(
\forall x\,(f+x=x)
\rightarrow
f=e
\right)
\right].
}
\]
\end{answerbox}

\end{workedexamplebox}

%%%%%%%%%%%%%%%%%%%%%%%%%%%%%%%%%%%%%%%%%%%%%%%%%%%%%%%%%%%%
\subsection{Bounded Quantifier Notation}
\label{subsec:bounded-quantifier-notation}
%%%%%%%%%%%%%%%%%%%%%%%%%%%%%%%%%%%%%%%%%%%%%%%%%%%%%%%%%%%%

Mathematics often writes

\[
\forall x\in A\,P(x).
\]

In ordinary first-order set-theoretic notation, this abbreviates

\[
\boxed{
\forall x
\,
(x\in A\rightarrow P(x)).
}
\]

Similarly,

\[
\exists x\in A\,P(x)
\]

abbreviates

\[
\boxed{
\exists x
\,
(x\in A\land P(x)).
}
\]

%%%%%%%%%%%%%%%%%%%%%%%%%%%%%%%%%%%%%%%%%%%%%%%%%%%%%%%%%%%%
\subsection{Logical Form of Definitions}
\label{subsec:logical-form-definitions}
%%%%%%%%%%%%%%%%%%%%%%%%%%%%%%%%%%%%%%%%%%%%%%%%%%%%%%%%%%%%

A mathematical definition often introduces a new predicate by an
equivalence.

For example, divisibility may be written

\[
\boxed{
a\mid b
\iff
\exists k\in\mathbb{Z}
\,
(b=ak).
}
\]

Formalization makes the logical structure explicit.

%%%%%%%%%%%%%%%%%%%%%%%%%%%%%%%%%%%%%%%%%%%%%%%%%%%%%%%%%%%%
\subsection{Logical Form of Theorems}
\label{subsec:logical-form-theorems}
%%%%%%%%%%%%%%%%%%%%%%%%%%%%%%%%%%%%%%%%%%%%%%%%%%%%%%%%%%%%

Many theorems have structure

\[
\boxed{
\forall x
\,
(
H(x)
\rightarrow
C(x)
).
}
\]

Here

\[
H(x)
\]

contains the hypotheses, while

\[
C(x)
\]

contains the conclusion.

A counterexample consists of an \(x\) satisfying

\[
\boxed{
H(x)\land\neg C(x).
}
\]

%%%%%%%%%%%%%%%%%%%%%%%%%%%%%%%%%%%%%%%%%%%%%%%%%%%%%%%%%%%%
\subsection{Formal Counterexamples}
\label{subsec:formal-counterexamples}
%%%%%%%%%%%%%%%%%%%%%%%%%%%%%%%%%%%%%%%%%%%%%%%%%%%%%%%%%%%%

To refute

\[
\forall x
\,
(H(x)\rightarrow C(x)),
\]

it suffices to find

\[
a
\]

such that

\[
\boxed{
H(a)
}
\]

and

\[
\boxed{
\neg C(a).
}
\]

This is exactly the quantified negation

\[
\boxed{
\exists x
\,
(H(x)\land\neg C(x)).
}
\]

%%%%%%%%%%%%%%%%%%%%%%%%%%%%%%%%%%%%%%%%%%%%%%%%%%%%%%%%%%%%
\subsection{Formal Logic and Mathematical Structures}
\label{subsec:logic-mathematical-structures}
%%%%%%%%%%%%%%%%%%%%%%%%%%%%%%%%%%%%%%%%%%%%%%%%%%%%%%%%%%%%

Formal logic separates:

\[
\boxed{
\text{the language describing structures}
}
\]

from

\[
\boxed{
\text{the structures themselves}.
}
\]

This makes it possible to compare very different mathematical objects that
satisfy the same axioms.

For example, many nonisomorphic groups satisfy the same first-order group
axioms.

%%%%%%%%%%%%%%%%%%%%%%%%%%%%%%%%%%%%%%%%%%%%%%%%%%%%%%%%%%%%
\subsection{Formal Logic and Computation}
\label{subsec:formal-logic-computation}
%%%%%%%%%%%%%%%%%%%%%%%%%%%%%%%%%%%%%%%%%%%%%%%%%%%%%%%%%%%%

Because formulas are finite symbolic objects, computers can manipulate
them algorithmically.

Examples include:

\[
\boxed{
\text{SAT solving},
}
\]

\[
\boxed{
\text{proof checking},
}
\]

\[
\boxed{
\text{term rewriting},
}
\]

and

\[
\boxed{
\text{automated theorem proving}.
}
\]

This creates a major bridge between logic and computer science.

%%%%%%%%%%%%%%%%%%%%%%%%%%%%%%%%%%%%%%%%%%%%%%%%%%%%%%%%%%%%
\subsection{Decision Problems}
\label{subsec:logical-decision-problems}
%%%%%%%%%%%%%%%%%%%%%%%%%%%%%%%%%%%%%%%%%%%%%%%%%%%%%%%%%%%%

A decision problem asks for a yes-or-no answer.

Examples include:

\[
\boxed{
\text{Is \(\varphi\) satisfiable?}
}
\]

and

\[
\boxed{
\text{Is \(\varphi\) valid?}
}
\]

For propositional logic, finite truth tables give a decision procedure.

For unrestricted first-order logic, the situation is fundamentally more
complicated.

%%%%%%%%%%%%%%%%%%%%%%%%%%%%%%%%%%%%%%%%%%%%%%%%%%%%%%%%%%%%
\subsection{Decidability Preview}
\label{subsec:formal-decidability-preview}
%%%%%%%%%%%%%%%%%%%%%%%%%%%%%%%%%%%%%%%%%%%%%%%%%%%%%%%%%%%%

A formal problem is decidable if an algorithm exists that always halts and
returns the correct answer.

Propositional validity is decidable.

General first-order validity is not decidable by an algorithm that always
halts on every input.

These computational limitations are studied more deeply in computability
theory.

%%%%%%%%%%%%%%%%%%%%%%%%%%%%%%%%%%%%%%%%%%%%%%%%%%%%%%%%%%%%
\subsection{Semi-Decidability Preview}
\label{subsec:formal-semidecidability-preview}
%%%%%%%%%%%%%%%%%%%%%%%%%%%%%%%%%%%%%%%%%%%%%%%%%%%%%%%%%%%%

Although unrestricted first-order validity is undecidable, valid
first-order formulas can be systematically enumerated through complete
proof systems.

Thus first-order validity has a form of effective recognizability.

This distinction between

\[
\boxed{
\text{decidable}
}
\]

and

\[
\boxed{
\text{recognizable}
}
\]

becomes central in computability theory.

%%%%%%%%%%%%%%%%%%%%%%%%%%%%%%%%%%%%%%%%%%%%%%%%%%%%%%%%%%%%
\subsection{Formal Logic and Foundations}
\label{subsec:formal-logic-foundations}
%%%%%%%%%%%%%%%%%%%%%%%%%%%%%%%%%%%%%%%%%%%%%%%%%%%%%%%%%%%%

Formal logic provides the language in which foundational questions can be
posed precisely.

Examples include:

\[
\boxed{
\text{What axioms are being assumed?}
}
\]

\[
\boxed{
\text{What can be formally proved from them?}
}
\]

\[
\boxed{
\text{Does the theory have a model?}
}
\]

\[
\boxed{
\text{Is the theory consistent?}
}
\]

and

\[
\boxed{
\text{What statements can the language express?}
}
\]

These questions motivate the remaining sections of Chapter~12.

%%%%%%%%%%%%%%%%%%%%%%%%%%%%%%%%%%%%%%%%%%%%%%%%%%%%%%%%%%%%
\subsection{Formal Logic Workflow}
\label{subsec:formal-logic-workflow}
%%%%%%%%%%%%%%%%%%%%%%%%%%%%%%%%%%%%%%%%%%%%%%%%%%%%%%%%%%%%

When analyzing a formal logical system:

\begin{enumerate}
    \item specify the alphabet or signature;
    \item define terms recursively;
    \item define formulas recursively;
    \item identify free and bound variables;
    \item specify the proof rules;
    \item define structures or truth assignments;
    \item define semantic satisfaction;
    \item distinguish \(\vdash\) from \(\models\);
    \item determine whether formulas are satisfiable or valid;
    \item translate informal mathematical statements carefully;
    \item track quantifier scope;
    \item avoid variable capture during substitution;
    \item identify the theory or set of assumptions;
    \item determine whether semantic consequence is intended;
    \item determine whether formal derivability is intended;
    \item use soundness to connect proof with truth;
    \item use completeness when semantic consequence must be converted into
          derivability;
    \item analyze models when studying consistency or expressibility;
    \item state clearly whether reasoning is propositional, first-order, or
          higher-order.
\end{enumerate}

%%%%%%%%%%%%%%%%%%%%%%%%%%%%%%%%%%%%%%%%%%%%%%%%%%%%%%%%%%%%
\subsection{Common Errors}
\label{subsec:formal-logic-common-errors}
%%%%%%%%%%%%%%%%%%%%%%%%%%%%%%%%%%%%%%%%%%%%%%%%%%%%%%%%%%%%

\subsubsection{Confusing \(\vdash\) and \(\models\)}

The expression

\[
\Gamma\vdash\varphi
\]

states formal provability.

The expression

\[
\Gamma\models\varphi
\]

states semantic consequence.

They are related by soundness and completeness but are not definitions of
the same concept.

%%%%%%%%%%%%%%%%%%%%%%%%%%%%%%%%%%%%%%%%%%%%%%%%%%%%%%%%%%%%
\subsubsection{Confusing a Formula With Its Meaning}

A formula is a syntactic object.

Its truth depends on an interpretation.

The same formal sentence may be true in one structure and false in
another.

%%%%%%%%%%%%%%%%%%%%%%%%%%%%%%%%%%%%%%%%%%%%%%%%%%%%%%%%%%%%
\subsubsection{Ignoring Quantifier Order}

In general,

\[
\boxed{
\forall x\exists y\,P(x,y)
}
\]

does not imply

\[
\boxed{
\exists y\forall x\,P(x,y).
}
\]

The first may use a different \(y\) for each \(x\).

%%%%%%%%%%%%%%%%%%%%%%%%%%%%%%%%%%%%%%%%%%%%%%%%%%%%%%%%%%%%
\subsubsection{Negating a Quantifier Without Changing Its Type}

The negation of

\[
\forall x\,P(x)
\]

is not

\[
\forall x\,\neg P(x).
\]

Instead,

\[
\boxed{
\neg\forall x\,P(x)
\equiv
\exists x\,\neg P(x).
}
\]

%%%%%%%%%%%%%%%%%%%%%%%%%%%%%%%%%%%%%%%%%%%%%%%%%%%%%%%%%%%%
\subsubsection{Variable Capture}

Substitution can accidentally place a free variable inside the scope of a
quantifier.

Bound variables should be renamed when necessary.

%%%%%%%%%%%%%%%%%%%%%%%%%%%%%%%%%%%%%%%%%%%%%%%%%%%%%%%%%%%%
\subsubsection{Treating Free Variables as If They Were Quantified}

The formula

\[
P(x)
\]

is not automatically a sentence.

Its truth may depend on the assignment to \(x\).

%%%%%%%%%%%%%%%%%%%%%%%%%%%%%%%%%%%%%%%%%%%%%%%%%%%%%%%%%%%%
\subsubsection{Confusing Satisfiable and Valid}

A satisfiable formula must be true under at least one interpretation.

A valid formula must be true under every interpretation.

Thus

\[
\boxed{
\text{valid}
\Longrightarrow
\text{satisfiable},
}
\]

but not conversely.

%%%%%%%%%%%%%%%%%%%%%%%%%%%%%%%%%%%%%%%%%%%%%%%%%%%%%%%%%%%%
\subsubsection{Confusing Contradiction With False Under One Assignment}

A contradiction is false under every interpretation.

A formula that happens to be false under one assignment may still be
satisfiable.

%%%%%%%%%%%%%%%%%%%%%%%%%%%%%%%%%%%%%%%%%%%%%%%%%%%%%%%%%%%%
\subsubsection{Assuming Every True Mathematical Statement Is Logically Valid}

For example,

\[
1+1=2
\]

is true in the intended arithmetic structure.

It is not a formula true in every arbitrary interpretation unless its
symbols and axioms are appropriately fixed.

Logical validity and truth in an intended model are different notions.

%%%%%%%%%%%%%%%%%%%%%%%%%%%%%%%%%%%%%%%%%%%%%%%%%%%%%%%%%%%%
\subsubsection{Ignoring the Domain}

The sentence

\[
\forall x\exists y\,(y>x)
\]

may be true in one ordered structure and false in another.

Truth always depends on the structure.

%%%%%%%%%%%%%%%%%%%%%%%%%%%%%%%%%%%%%%%%%%%%%%%%%%%%%%%%%%%%
\subsubsection{Confusing Logical Completeness With Theory Completeness}

A proof system is complete when every semantic consequence is formally
derivable.

A theory is complete when, for every sentence \(\varphi\), it decides
\(\varphi\) or its negation.

These are different concepts.

%%%%%%%%%%%%%%%%%%%%%%%%%%%%%%%%%%%%%%%%%%%%%%%%%%%%%%%%%%%%
\subsubsection{Confusing Gödel Completeness and Incompleteness}

Gödel's completeness theorem states that first-order semantic consequence
matches formal derivability.

Gödel's incompleteness theorems state that sufficiently strong consistent
effective theories cannot decide every relevant arithmetic statement.

There is no contradiction between these results.

%%%%%%%%%%%%%%%%%%%%%%%%%%%%%%%%%%%%%%%%%%%%%%%%%%%%%%%%%%%%
\subsubsection{Assuming First-Order Logic Can Express Every Mathematical Property}

First-order quantifiers range only over domain elements.

Properties requiring unrestricted quantification over subsets or relations
may require stronger languages.

%%%%%%%%%%%%%%%%%%%%%%%%%%%%%%%%%%%%%%%%%%%%%%%%%%%%%%%%%%%%
\subsection{Applications}
\label{subsec:formal-logic-applications}
%%%%%%%%%%%%%%%%%%%%%%%%%%%%%%%%%%%%%%%%%%%%%%%%%%%%%%%%%%%%

\begin{applicationbox}

\textbf{Foundations of Mathematics.}
Formal logic provides precise languages and proof systems for axiomatic
mathematics.

\medskip

\textbf{Set Theory.}
Set-theoretic axioms are expressed in a formal first-order language.

\medskip

\textbf{Abstract Algebra.}
Groups, rings, fields, and other structures can be studied as models of
formal theories.

\medskip

\textbf{Computer Science.}
Logic underlies program verification, SAT solving, databases, formal
methods, type systems, and automated reasoning.

\medskip

\textbf{Model Theory.}
Model theory studies relationships between formal theories and their
mathematical structures.

\medskip

\textbf{Proof Theory.}
Proof theory studies formal derivations as mathematical objects.

\medskip

\textbf{Computability Theory.}
Logical languages allow algorithmic solvability and undecidability to be
defined precisely.

\medskip

\textbf{Automated Theorem Proving.}
Formal inference rules can be implemented computationally to search for
proofs and verify deductions.

\end{applicationbox}

%%%%%%%%%%%%%%%%%%%%%%%%%%%%%%%%%%%%%%%%%%%%%%%%%%%%%%%%%%%%
\subsection{Exercises}
\label{subsec:formal-logic-exercises}
%%%%%%%%%%%%%%%%%%%%%%%%%%%%%%%%%%%%%%%%%%%%%%%%%%%%%%%%%%%%

\begin{exercise}[1]
Define syntax.
\end{exercise}

\begin{exercise}[1]
Define semantics.
\end{exercise}

\begin{exercise}[1]
Distinguish object language and metalanguage.
\end{exercise}

\begin{exercise}[1]
Define a well-formed formula.
\end{exercise}

\begin{exercise}[1]
Define satisfiability.
\end{exercise}

\begin{exercise}[1]
Define validity.
\end{exercise}

\begin{exercise}[1]
Define semantic consequence.
\end{exercise}

\begin{exercise}[1]
Define syntactic derivability.
\end{exercise}

\begin{exercise}[1]
Define soundness.
\end{exercise}

\begin{exercise}[1]
Define completeness.
\end{exercise}

\begin{exercise}[2]
Construct the parse structure of

\[
\neg(P\lor(Q\land R)).
\]
\end{exercise}

\begin{exercise}[2]
List all subformulas of

\[
(P\rightarrow Q)\land\neg R.
\]
\end{exercise}

\begin{exercise}[2]
Determine whether

\[
P\lor Q
\]

is satisfiable.
\end{exercise}

\begin{exercise}[2]
Determine whether

\[
P\land\neg P
\]

is satisfiable.
\end{exercise}

\begin{exercise}[2]
Show using a truth table that

\[
P\rightarrow Q
\equiv
\neg P\lor Q.
\]
\end{exercise}

\begin{exercise}[2]
Verify De Morgan's laws using truth assignments.
\end{exercise}

\begin{exercise}[2]
Show that

\[
P\rightarrow Q
\]

is equivalent to its contrapositive.
\end{exercise}

\begin{exercise}[3]
Convert

\[
P\rightarrow(Q\land R)
\]

to CNF.
\end{exercise}

\begin{exercise}[3]
Convert

\[
(P\land Q)\rightarrow R
\]

to CNF.
\end{exercise}

\begin{exercise}[3]
Construct a DNF for a given three-variable truth table.
\end{exercise}

\begin{exercise}[3]
Show that

\[
\{\neg,\land\}
\]

is functionally complete.
\end{exercise}

\begin{exercise}[3]
Express negation using only NAND.
\end{exercise}

\begin{exercise}[3]
Express conjunction using only NAND.
\end{exercise}

\begin{exercise}[3]
Determine how many truth assignments exist for a formula with \(12\)
variables.
\end{exercise}

\begin{exercise}[2]
Explain the difference between

\[
\Gamma\vdash\varphi
\]

and

\[
\Gamma\models\varphi.
\]
\end{exercise}

\begin{exercise}[2]
Use modus ponens to derive \(Q\) from

\[
P
\]

and

\[
P\rightarrow Q.
\]
\end{exercise}

\begin{exercise}[3]
Construct a natural-deduction proof of

\[
P\land Q
\rightarrow
P.
\]
\end{exercise}

\begin{exercise}[3]
Construct a natural-deduction proof of

\[
P
\rightarrow
(Q\rightarrow P).
\]
\end{exercise}

\begin{exercise}[3]
Explain implication introduction.
\end{exercise}

\begin{exercise}[3]
Explain proof by contradiction in a formal proof system.
\end{exercise}

\begin{exercise}[3]
Explain the principle of explosion.
\end{exercise}

\begin{exercise}[2]
In a first-order language of arithmetic, identify which of the following
are terms and which are formulas:

\[
x+1,
\qquad
x<y,
\qquad
(x+y)\cdot z,
\qquad
x=y.
\]
\end{exercise}

\begin{exercise}[2]
Identify all free and bound variables in

\[
\forall x
\,
(P(x,y)\rightarrow\exists z\,Q(x,z)).
\]
\end{exercise}

\begin{exercise}[3]
Rename bound variables in a formula without changing its meaning.
\end{exercise}

\begin{exercise}[3]
Give an example of variable capture.
\end{exercise}

\begin{exercise}[3]
Perform a capture-avoiding substitution.
\end{exercise}

\begin{exercise}[2]
Determine whether

\[
\forall x\exists y\,(x<y)
\]

is true over \(\mathbb{R}\).
\end{exercise}

\begin{exercise}[2]
Determine whether

\[
\exists y\forall x\,(x<y)
\]

is true over \(\mathbb{R}\).
\end{exercise}

\begin{exercise}[3]
Negate

\[
\forall x\exists y\,P(x,y)
\]

and move the negation completely inward.
\end{exercise}

\begin{exercise}[3]
Negate

\[
\exists x
\forall y
\,
(P(x,y)\rightarrow Q(y)).
\]
\end{exercise}

\begin{exercise}[2]
Formalize:

\[
\text{``Every integer has an additive inverse.''}
\]
\end{exercise}

\begin{exercise}[2]
Formalize:

\[
\text{``There exists a largest real number.''}
\]
\end{exercise}

\begin{exercise}[3]
Formalize:

\[
\text{``Every positive real number has a positive square root.''}
\]
\end{exercise}

\begin{exercise}[3]
Expand

\[
\exists!x\,P(x)
\]

using only \(\exists,\forall,=,\land,\rightarrow\).
\end{exercise}

\begin{exercise}[3]
Rewrite

\[
\forall x\in A\,P(x)
\]

using unrestricted quantification and membership.
\end{exercise}

\begin{exercise}[3]
Rewrite

\[
\exists x\in A\,P(x)
\]

using unrestricted quantification and membership.
\end{exercise}

\begin{exercise}[2]
Define a first-order structure for a language containing one binary
relation symbol.
\end{exercise}

\begin{exercise}[3]
Interpret the language

\[
\{0,+,<\}
\]

over the integers.
\end{exercise}

\begin{exercise}[3]
Evaluate a term under a supplied structure and variable assignment.
\end{exercise}

\begin{exercise}[3]
Determine whether a supplied structure satisfies a first-order sentence.
\end{exercise}

\begin{exercise}[3]
Write first-order axioms for a commutative binary operation.
\end{exercise}

\begin{exercise}[3]
Write first-order axioms for a partial order.
\end{exercise}

\begin{exercise}[3]
Write first-order axioms for an equivalence relation.
\end{exercise}

\begin{exercise}[3]
Write the associativity axiom for a group language.
\end{exercise}

\begin{exercise}[3]
Formalize existence and uniqueness of a group identity.
\end{exercise}

\begin{exercise}[3]
Explain why a theory is a set of sentences rather than merely one formula.
\end{exercise}

\begin{exercise}[3]
Explain the relationship between a theory and its models.
\end{exercise}

\begin{exercise}[3]
Explain the difference between truth in one structure and logical
validity.
\end{exercise}

\begin{exercise}[3]
Give an example of a satisfiable but nonvalid first-order sentence.
\end{exercise}

\begin{exercise}[3]
Explain universal elimination.
\end{exercise}

\begin{exercise}[3]
Explain the side condition required for universal introduction.
\end{exercise}

\begin{exercise}[3]
Explain existential introduction.
\end{exercise}

\begin{exercise}[3]
Explain why existential elimination requires a fresh representative.
\end{exercise}

\begin{exercise}[3]
State first-order soundness.
\end{exercise}

\begin{exercise}[3]
State first-order completeness.
\end{exercise}

\begin{exercise}[3]
Explain why soundness and completeness connect syntax and semantics.
\end{exercise}

\begin{exercise}[3]
Distinguish logical completeness from completeness of a theory.
\end{exercise}

\begin{exercise}[3]
Explain the compactness theorem conceptually.
\end{exercise}

\begin{exercise}[3]
Use compactness to explain why a theory containing sentences asserting
``there are at least \(n\) elements'' for every \(n\) has an infinite
model.
\end{exercise}

\begin{exercise}[3]
Explain the Löwenheim--Skolem phenomenon conceptually.
\end{exercise}

\begin{exercise}[3]
Explain why first-order logic does not quantify directly over arbitrary
subsets of the domain.
\end{exercise}

\begin{exercise}[3]
Distinguish first-order and second-order quantification.
\end{exercise}

\begin{exercise}[3]
Express the statement

\[
\text{``not every real number is positive''}
\]

using quantifiers and then negate it correctly.
\end{exercise}

\begin{exercise}[3]
Formalize a theorem of the form

\[
\text{``Every object satisfying \(H\) also satisfies \(C\).''}
\]
\end{exercise}

\begin{exercise}[3]
Formalize the logical form of a counterexample to that theorem.
\end{exercise}

\begin{exercise}[3]
Explain why propositional validity is decidable.
\end{exercise}

\begin{exercise}[3]
Explain why first-order logic raises deeper computability questions than
propositional logic.
\end{exercise}

\begin{exercise}[4]
Prove unique readability for a simple propositional grammar.
\end{exercise}

\begin{exercise}[4]
Prove structural induction for propositional formulas.
\end{exercise}

\begin{exercise}[4]
Prove that every propositional formula has an equivalent CNF.
\end{exercise}

\begin{exercise}[4]
Prove that every propositional formula has an equivalent DNF.
\end{exercise}

\begin{exercise}[4]
Study functional completeness of NAND.
\end{exercise}

\begin{exercise}[4]
Study functional completeness of NOR.
\end{exercise}

\begin{exercise}[4]
Study the semantics of natural deduction.
\end{exercise}

\begin{exercise}[4]
Prove soundness of selected propositional inference rules.
\end{exercise}

\begin{exercise}[4]
Study completeness of propositional logic.
\end{exercise}

\begin{exercise}[4]
Study proof systems based on Hilbert calculi.
\end{exercise}

\begin{exercise}[4]
Study sequent calculus.
\end{exercise}

\begin{exercise}[4]
Study the deduction theorem.
\end{exercise}

\begin{exercise}[4]
Prove the substitution lemma for first-order terms.
\end{exercise}

\begin{exercise}[4]
Prove the substitution lemma for first-order formulas.
\end{exercise}

\begin{exercise}[4]
Study alpha-equivalence of formulas.
\end{exercise}

\begin{exercise}[4]
Study prenex normal form.
\end{exercise}

\begin{exercise}[4]
Study Skolemization.
\end{exercise}

\begin{exercise}[4]
Study Herbrand-style methods.
\end{exercise}

\begin{exercise}[4]
Study resolution for propositional logic.
\end{exercise}

\begin{exercise}[4]
Study first-order resolution.
\end{exercise}

\begin{exercise}[4]
Study unification algorithms.
\end{exercise}

\begin{exercise}[4]
Study the proof of first-order completeness.
\end{exercise}

\begin{exercise}[4]
Study the compactness theorem and its consequences.
\end{exercise}

\begin{exercise}[4]
Study downward Löwenheim--Skolem.
\end{exercise}

\begin{exercise}[4]
Study upward Löwenheim--Skolem.
\end{exercise}

\begin{exercise}[4]
Investigate the boundary between first-order and second-order
expressibility.
\end{exercise}

%%%%%%%%%%%%%%%%%%%%%%%%%%%%%%%%%%%%%%%%%%%%%%%%%%%%%%%%%%%%
\subsection{Conceptual Summary}
\label{subsec:formal-logic-summary}
%%%%%%%%%%%%%%%%%%%%%%%%%%%%%%%%%%%%%%%%%%%%%%%%%%%%%%%%%%%%

Formal logic treats mathematical language and deduction as mathematical
objects.

Its first fundamental distinction is

\[
\boxed{
\text{syntax}
\quad\text{versus}\quad
\text{semantics}.
}
\]

Syntactic derivability is written

\[
\boxed{
\Gamma\vdash\varphi,
}
\]

while semantic consequence is written

\[
\boxed{
\Gamma\models\varphi.
}
\]

Soundness gives

\[
\boxed{
\Gamma\vdash\varphi
\Longrightarrow
\Gamma\models\varphi,
}
\]

while completeness gives

\[
\boxed{
\Gamma\models\varphi
\Longrightarrow
\Gamma\vdash\varphi.
}
\]

Propositional logic studies formulas built from truth-valued variables and
logical connectives.

First-order logic adds:

\[
\boxed{
\text{objects}
+
\text{functions}
+
\text{relations}
+
\text{quantifiers}.
}
\]

A first-order structure gives semantic meaning to a formal language.

A theory is a set of sentences, and its models are the structures satisfying
those sentences.

Formal logic therefore creates the framework for studying

\[
\boxed{
\text{proof}
+
\text{truth}
+
\text{models}
+
\text{expressibility}
+
\text{computation}.
}
\]

%%%%%%%%%%%%%%%%%%%%%%%%%%%%%%%%%%%%%%%%%%%%%%%%%%%%%%%%%%%%
\subsection{Section Connections}
\label{subsec:formal-logic-connections}
%%%%%%%%%%%%%%%%%%%%%%%%%%%%%%%%%%%%%%%%%%%%%%%%%%%%%%%%%%%%

\chapterconnection{
Formal Logic begins Chapter~12 by converting the elementary logical
reasoning introduced in Chapter~0 into a precise mathematical theory of
syntax, semantics, structures, and formal proofs. Propositional logic
introduces satisfiability, validity, semantic consequence, derivability,
soundness, and completeness, while first-order logic adds terms,
quantifiers, mathematical structures, theories, and models. These ideas
provide the formal language required for the next section, Axiomatic Set
Theory, where sets and membership are governed by explicit axioms. Formal
Logic also supplies the foundation for the later sections on Model Theory,
Proof Theory, Computability Theory, Recursion Theory, Type Theory,
Constructive Mathematics, and the Philosophy and Foundations of
Mathematics.
}

%%%%%%%%%%%%%%%%%%%%%%%%%%%%%%%%%%%%%%%%%%%%%%%%%%%%%%%%%%%%
% SECTION SOURCE TRACKING
%%%%%%%%%%%%%%%%%%%%%%%%%%%%%%%%%%%%%%%%%%%%%%%%%%%%%%%%%%%%

%------------------------------------------------------------
% 12.1 Formal Logic
%------------------------------------------------------------
%
% Source basis:
%   - Standard classical propositional logic.
%   - Standard first-order logic.
%   - Standard proof-theoretic and model-theoretic foundations.
%
% Used for:
%   - syntax and semantics;
%   - object language and metalanguage;
%   - formal languages;
%   - recursive syntax;
%   - propositional logic;
%   - well-formed formulas;
%   - parse trees and unique readability;
%   - truth assignments;
%   - satisfaction;
%   - satisfiability and validity;
%   - contradictions;
%   - logical consequence;
%   - logical equivalence;
%   - propositional equivalences;
%   - CNF and DNF;
%   - Boolean functions;
%   - functional completeness;
%   - NAND and NOR;
%   - SAT;
%   - formal proof systems;
%   - syntactic derivability;
%   - inference rules;
%   - natural deduction;
%   - implication introduction;
%   - contradiction;
%   - deduction theorem;
%   - soundness;
%   - completeness;
%   - consistency;
%   - principle of explosion;
%   - first-order signatures;
%   - variables, constants, functions, and relations;
%   - terms;
%   - atomic formulas;
%   - first-order formulas;
%   - free and bound variables;
%   - sentences;
%   - quantifier scope;
%   - substitution;
%   - variable capture;
%   - bound-variable renaming;
%   - structures and interpretations;
%   - variable assignments;
%   - term evaluation;
%   - first-order satisfaction;
%   - universal and existential semantics;
%   - quantifier negation;
%   - models and theories;
%   - first-order validity and satisfiability;
%   - semantic consequence;
%   - group axioms as first-order theories;
%   - equality;
%   - quantifier inference rules;
%   - first-order soundness and completeness;
%   - consistency and satisfiability;
%   - compactness preview;
%   - Löwenheim--Skolem preview;
%   - expressive power and limitations;
%   - first- versus second-order logic;
%   - translation of mathematical statements;
%   - uniqueness quantifier;
%   - bounded-quantifier notation;
%   - logical form of definitions and theorems;
%   - formal counterexamples;
%   - decision problems;
%   - decidability and recognizability previews;
%   - applications and exercises.
%
% Notes:
%   - Axiomatic Set Theory is developed next in Section 12.2.
%   - Model Theory follows in Section 12.3.
%   - Proof Theory follows in Section 12.4.
%   - Computability Theory follows in Section 12.5.
%   - Elementary Logic appears in Chapter 0.
%
%%%%%%%%%%%%%%%%%%%%%%%%%%%%%%%%%%%%%%%%%%%%%%%%%%%%%%%%%%%%